
\documentclass[12pt,a4paper]{report}

% Packages
\usepackage[a4paper, top=25mm, bottom=30mm, left=22mm, right=22mm, headheight=18pt, headsep=12pt, footskip=16mm]{geometry}
\usepackage{graphicx}
\usepackage{amsmath}
\usepackage{amssymb}
\usepackage{booktabs}
\usepackage{array}
\usepackage{tabularx}
\usepackage{float}
\usepackage{fancyhdr}
\usepackage[hidelinks]{hyperref}
\usepackage{xcolor}
\usepackage{setspace}
\usepackage{enumitem}
\usepackage{caption}

\captionsetup{
    labelfont=bf,
    justification=raggedright,
    singlelinecheck=false,
    format=plain
}
\usepackage{subcaption}
\usepackage{multirow}
\usepackage{colortbl}
\usepackage{longtable}
\usepackage{titlesec}
\usepackage{titletoc}
\usepackage{lastpage}
\usepackage{makecell}
\usepackage{etoolbox}
\usepackage{needspace}

\numberwithin{table}{chapter}
\numberwithin{figure}{chapter}


% Table layout and spacing
\setlength{\tabcolsep}{5pt}
\renewcommand{\arraystretch}{1.15}
\setlength{\extrarowheight}{1.2pt}
\setlength{\arrayrulewidth}{0.5pt}
\setlength{\LTpre}{0pt}
\setlength{\LTpost}{6pt}
\setlength{\LTleft}{\fill}
\setlength{\LTright}{\fill}
\BeforeBeginEnvironment{table}{\needspace{4\baselineskip}}
\BeforeBeginEnvironment{longtable}{\needspace{4\baselineskip}}


\definecolor{osdagGreen}{HTML}{91B014}
\definecolor{osdagGreenDark}{HTML}{6E870D}
\definecolor{osdagSlate}{HTML}{2C3E50}
\definecolor{osdagLightGray}{HTML}{F4F6F8}
\definecolor{osdagBorder}{HTML}{D0D0D0}
\definecolor{osdagPass}{HTML}{2E7D32}
\definecolor{osdagWarn}{HTML}{E65100}
\definecolor{osdagFail}{HTML}{C62828}
\definecolor{osdagHeaderBg}{HTML}{E8EFA8}


\fancypagestyle{main}{
  \fancyhf{}
  \fancyhead[L]{\small\textcolor{osdagSlate}{Four-Lane Highway Bridge over River Kali $|$ NHAI-KA-2026-BR-04}}
  \fancyhead[R]{\small\textcolor{osdagSlate}{2026-11-20 $|$ Rev 0}}
  \fancyfoot[L]{\small\textcolor{osdagSlate}{Osdag $|$ FOSSEE $|$ Indian Institute of Technology Bombay}}
  \fancyfoot[R]{\small\textcolor{osdagSlate}{Page \thepage\ of \pageref{LastPage}}}
  \renewcommand{\headrule}{\color{osdagGreen}\hrule width\headwidth height 1pt \vspace{2pt}}
  \renewcommand{\footrule}{%
    \ifbool{hasSDonPage}{%
      \vspace{-12pt}%
      \hbox to \headwidth{\textcolor{black}{\footnotesize\textit{* Software default value}}\hfil}%
      \vspace{2pt}%
    }{%
      \vspace{-4pt}%
    }%
    \color{osdagGreen}\hrule width\headwidth height 1pt \vspace{4pt}%
  }
}

\fancypagestyle{plain}{
  \fancyhf{}
  \fancyhead[L]{\small\textcolor{osdagSlate}{Four-Lane Highway Bridge over River Kali $|$ NHAI-KA-2026-BR-04}}
  \fancyhead[R]{\small\textcolor{osdagSlate}{2026-11-20 $|$ Rev 0}}
  \fancyfoot[L]{\small\textcolor{osdagSlate}{Osdag $|$ FOSSEE $|$ Indian Institute of Technology Bombay}}
  \fancyfoot[R]{\small\textcolor{osdagSlate}{Page \thepage\ of \pageref{LastPage}}}
  \renewcommand{\headrule}{\color{osdagGreen}\hrule width\headwidth height 1pt \vspace{2pt}}
  \renewcommand{\footrule}{%
    \ifbool{hasSDonPage}{%
      \vspace{-12pt}%
      \hbox to \headwidth{\textcolor{black}{\footnotesize\textit{* Software default value}}\hfil}%
      \vspace{2pt}%
    }{%
      \vspace{-4pt}%
    }%
    \color{osdagGreen}\hrule width\headwidth height 1pt \vspace{4pt}%
  }
}

\fancypagestyle{firstpage}{
  \fancyhf{}
  \renewcommand{\headrulewidth}{0pt}
  \fancyfoot[L]{\small\textcolor{osdagSlate}{Osdag $|$ FOSSEE $|$ Indian Institute of Technology Bombay}}
  \fancyfoot[R]{\small\textcolor{osdagSlate}{Page \thepage\ of \pageref{LastPage}}}
  \renewcommand{\footrule}{\color{osdagGreen}\hrule width\headwidth height 1pt \vspace{4pt}}
}

\pagestyle{main}
\setstretch{1.15}

% Custom Commands
\newcommand{\placeholder}[1]{\textit{\textless #1\textgreater}}
\newcommand{\todo}[1]{\colorbox{yellow}{TODO: #1}}
\newcolumntype{L}[1]{>{\raggedright\arraybackslash}p{#1}}
\newcolumntype{C}[1]{>{\centering\arraybackslash}p{#1}}
\newcolumntype{R}[1]{>{\raggedleft\arraybackslash}p{#1}}

% Software-default asterisk
\newcommand{\sdstar}{\textsuperscript{*}}
\newbool{hasSDonPage}
\boolfalse{hasSDonPage}
\newcommand{\markSD}{\global\booltrue{hasSDonPage}}
\renewcommand{\sdstar}{\textsuperscript{*}\markSD{}}
\AddToHook{shipout/before}{\global\boolfalse{hasSDonPage}}

\title{\Large\textbf{OsdagBridge} \\ \normalsize Open Source Software for Steel Girder Bridge Design \\ \vspace{2cm} \large Design Report}
\author{}
\date{}

\begin{document}


\begin{titlepage}
\thispagestyle{firstpage}
\centering
\vspace*{1.5cm}
\noindent
\begin{minipage}[c]{0.6\textwidth}
\raggedright
\includegraphics[width=\linewidth,keepaspectratio]{assets/osdag_logo.png}
\end{minipage}%
\hfill
\begin{minipage}[c]{0.35\textwidth}
\raggedleft

\end{minipage}
\\[1cm]

{\Huge \textbf{OsdagBridge}}\\[0.3cm]
{\large Open Source Software for Steel Girder Bridge Design}\\[1.5cm]
{\Large Design Report}\\[1.5cm]
\begin{tabular}{|L{4cm}|L{10cm}|}
\hline
\textbf{Project Name} & Four-Lane Highway Bridge over River Kali \\
\hline
\textbf{Project Location} & Karwar, Karnataka \\
\hline
\textbf{Author / Designer} & Er. Rajesh Sharma \\
\hline
\textbf{Reviewer} & Dr. S. K. Bhattacharya \\
\hline
\textbf{Organization} & FOSSEE Structural Engineering Group \\
\hline
\textbf{Client Name and Organization} & National Highways Authority of India (NHAI) \\
\hline
\textbf{Job Number} & NHAI-KA-2026-BR-04 \\
\hline
\textbf{Date} & 2026-11-20 \\
\hline
\textbf{Report Version} & Rev 0 \\
\hline
\end{tabular}
\end{titlepage}


% Chapter / TOC Formatting
\titleformat{\chapter}[block]
  {\normalfont\Large\bfseries\centering}{\thechapter}{1em}{}
\titlespacing*{\chapter}{0pt}{0pt}{10pt}
\setcounter{tocdepth}{2}

% TOC styling using titletoc
\titlecontents{chapter}[1.5em]
  {\normalfont\vspace{2pt}}
  {\contentslabel{1.5em}}
  {\hspace*{-1.5em}}
  {\hfill\contentspage}

\titlecontents{section}[3.8em]
  {\normalfont}
  {\contentslabel{2.3em}}
  {\hspace*{-2.3em}}
  {\hfill\contentspage}

\titlecontents{subsection}[7.0em]
  {\normalfont}
  {\contentslabel{3.2em}}
  {\hspace*{-3.2em}}
  {\hfill\contentspage}

\newpage
\renewcommand{\contentsname}{\centering\Large\bfseries Table of Contents}
\tableofcontents


\newpage
{\centering\Large\bfseries Executive Summary\par}
\addcontentsline{toc}{chapter}{Executive Summary}
\vspace{0.8em}

This section provides a concise summary of the bridge design, key inputs, governing loads, and final design outcomes.

\section*{Project Overview}
\addcontentsline{toc}{section}{Project Overview}
\label{sec:project-overview}


\begin{tabular}{|L{5.5cm}|L{8.5cm}|}
\hline
\textbf{Bridge Type} & Plate Girder Bridge \\
\hline
\textbf{Design Standard} & IRC 5, IRC 6, IRC 22, IRC 24, IS 800 \\
\hline
\textbf{Span} & 30.0 m \\
\hline
\textbf{Carriageway Width} & 7.5 m \\
\hline
\textbf{No. of Girders} & 4 \\
\hline
\textbf{Girder Spacing} & 2.5 \\
\hline
\textbf{Deck Thickness} & 250 \\
\hline
\textbf{Overall Design Status} & Pass \\
\hline
\textbf{Governing Check} & Flexure \\
\hline
\textbf{Overall Utilization Ratio (max)} & 0.82 (Deck slab) \\
\hline
\end{tabular}


\begin{figure}[H]
\centering
\fbox{\parbox{0.97\textwidth}{
\textit{[ PLACEHOLDER: Figure 1 -- Overall Bridge Plan ]}
}}
\caption*{Figure 1 -- Overall Bridge Plan}
\end{figure}

\newpage

\begin{figure}[H]
\centering
\fbox{\parbox{0.97\textwidth}{
\textit{[ PLACEHOLDER: Figure 2 -- Typical Cross-Section (with girder, deck, barriers, footpath) ]}
}}
\caption*{Figure 2 -- Typical Cross-Section (with girder, deck, barriers, footpath)}
\end{figure}

\begin{figure}[H]
\centering
\fbox{\parbox{0.97\textwidth}{
\textit{[ PLACEHOLDER: Figure 3 -- 3D View of Bridge Superstructure ]}
}}
\caption*{Figure 3 -- 3D View of Bridge Superstructure}
\end{figure}

\noindent\textbf{Table 1 -- Final Bridge Geometry (after optimization)}

\vspace{0.4em}
\noindent
\begin{tabular}{|C{2.8cm}|C{3.0cm}|C{3.0cm}|C{3.0cm}|C{3.0cm}|}
\hline
  \textbf{} &
  \textbf{G1} &
  \textbf{G2} &
  \textbf{G3} &
  \textbf{G4} \\
\hline
Member ID & G1M1 & G2M1 & G3M1 & G4M1 \\
\hline
Section Designation & PG-1864x550 & PG-1864x550 & PG-1864x550 & PG-1864x550 \\
\hline
Governing Check & Flexure & Flexure & Flexure & Flexure \\
\hline
Utilization Ratio & 0.82 (Deck slab) & 0.82 (Deck slab) & 0.82 (Deck slab) & 0.82 (Deck slab) \\
\hline
\end{tabular}

\vspace{0.4em}
\noindent\textit{Note: Utilization ratio (UR) = demand / capacity. A value $< 1.0$ indicates a passing check.}

\vspace{1em}

\section*{Key Design Outcomes Summary}
\addcontentsline{toc}{section}{Key Design Outcomes Summary}
\label{sec:key-outcomes}

\noindent Girder design pass \\
Cross bracing design pass \\
End Diaphragm design pass \\
Deck design pass

\section*{Design Assumptions and Limitations}
\addcontentsline{toc}{section}{Design Assumptions and Limitations}
\label{sec:assumptions}

\begin{itemize}
\item Additional inputs not provided by the user were assumed by software per IRC/IS code defaults or practical consideration.
\item Grillage analysis was performed using OSPGrillage assuming simply supported I-girders.
\item Substructure and foundation design are not included in this report.
\item Splice connections and bearings are not designed in this version.
\end{itemize}

% Restore numbered chapter format
\titleformat{\chapter}[block]{\normalfont\Large\bfseries\centering}{\thechapter}{1em}{}
\titlespacing*{\chapter}{0pt}{-30pt}{10pt}


\chapter{Project Information}

This section records all project metadata as entered by the designer.

\section{Project and Design Team Details}
\label{sec:project-details}

\begin{tabular}{|L{5.5cm}|L{8.5cm}|}
\hline
\textbf{Project Name} & Four-Lane Highway Bridge over River Kali \\
\hline
\textbf{Project Location} & Karwar, Karnataka \\
\hline
\textbf{Designer} & Er. Rajesh Sharma \\
\hline
\textbf{Reviewer} & Dr. S. K. Bhattacharya \\
\hline
\textbf{Organization} & FOSSEE Structural Engineering Group \\
\hline
\textbf{Client} & National Highways Authority of India (NHAI) \\
\hline
\textbf{Software Version} & OsdagBridge \\
\hline
\end{tabular}


\section{Applicable Codes and Standards}
\label{sec:codes}

\begin{itemize}
\item Indian Roads Congress (IRC) 5: General Features of Design
\item Indian Roads Congress (IRC) 6: Loads and Load Combinations
\item Indian Roads Congress (IRC) 22: Composite Construction (Limit State Design)
\item Indian Roads Congress (IRC) 24: Steel Road Bridges (Limit State Method)
\item Indian Roads Congress (IRC) 112: Concrete Road Bridges (deck design)
\item Indian Roads Congress Special Publication (IRC SP) 114: Seismic Design of Road Bridges
\item Indian Standard (IS) 800: General Construction in Steel
\item Indian Standard (IS) 2062: Hot Rolled Structural Steel Specification
\item Indian Standard (IS) 6006: Steel Bearings
\end{itemize}


\chapter{Input Parameters}

\setlength{\abovecaptionskip}{2pt}
\setlength{\belowcaptionskip}{2pt}

This section documents all inputs provided to OsdagBridge. User-provided inputs are clearly distinguished from software-assumed defaults. Where the user did not supply a value, the software has applied the IRC/IS code default or an empirical guideline; these are annotated with an asterisk (\sdstar{}).

\section{Basic Inputs (User-Defined)}
\label{sec:basic-inputs}

\noindent\textit{Note: These inputs are mandatory and were provided by the user.}

\begin{table}[H]
\caption{\textbf{Project Location}}
\label{subsec:project-location}
\begin{tabular}{|L{5.5cm}|L{8.5cm}|}
\hline
\textbf{parameter} & \textbf{value} \\
\hline
\textnormal{Project Location} & Karwar, Karnataka \\
\hline
\textnormal{Latitude / Longitude} & 19.0760, 72.8777 \\
\hline
\textnormal{Seismic Zone (IRC 6)} & Zone IV \\
\hline
\textnormal{Basic Wind Speed (IRC 6)} & 39.0 m/s \\
\hline
\textnormal{Shade Temp. Max / Min (IRC 6)} & 45.0 °C / -5.0 °C \\
\hline
\end{tabular}
\vspace{0.4cm}
\end{table}

\begin{table}[H]
\caption{\textbf{Bridge Geometry}}
\label{subsec:bridge-geometry}
\begin{tabular}{|L{5.5cm}|L{8.5cm}|}
\hline
\textbf{parameter} & \textbf{value} \\
\hline
\textnormal{Type of Structure} & Plate Girder Bridge \\
\hline
\textnormal{Span (m)} & 30.0 m \\
\hline
\textnormal{Carriageway Width (m)} & 7.5 m \\
\hline
\textnormal{Include Median} & No \\
\hline
\textnormal{Footpath} & Both Sides \\
\hline
\textnormal{Skew Angle (degrees)} & 0° (IRC 24 Cl. 504.8 limit: $\pm$15°) \\
\hline
\end{tabular}
\end{table}
\vspace{0.4cm}

\begin{table}[H]
\caption{\textbf{Material Selection}}
\label{subsec:material}
\begin{tabular}{|L{5.5cm}|L{8.5cm}|}
\hline
\textbf{parameter} & \textbf{value} \\
\hline
\textnormal{Girder Steel Grade (IS 2062)} & E 250 (Fe 410 W) B \\
\hline
\textnormal{Cross Bracing Steel Grade} & E 250 (Fe 410 W) B \\
\hline
\textnormal{End Diaphragm Steel Grade} & E 250 (Fe 410 W) B \\
\hline
\textnormal{Concrete Deck Grade (IRC 22)} & M35 \\
\hline
\end{tabular}
\end{table}
\vspace{0.4cm}

\newpage
\section{Additional Inputs}
\label{sec:additional-inputs}

Where the user has modified additional inputs, those values are reported here. Where no modification was made, the software default is shown.

\vspace{0.8cm}

\begin{longtable}{|L{5.5cm}|p{10.0cm}|}
\caption{\textbf{Typical Section Details}}\label{tab:typical-section}\\
\hline
\textbf{Parameter} & \textbf{Value} \\
\hline
\endfirsthead

\multicolumn{2}{c}{\small\textit{\tablename\ \thetable\ -- Continued from previous page}} \\
\hline
\textbf{Parameter} & \textbf{Value} \\
\hline
\endhead

\hline
\multicolumn{2}{r}{\footnotesize\textit{Continued on next page\ldots}} \\
\endfoot

\hline
\endlastfoot

\textnormal{Overall Bridge Width (m)} & 10.5 \\[6pt]
\hline
\textnormal{No. of Girders} & 4 \\[6pt]
\hline
\textnormal{Girder Spacing (m)} & 2.5 m \\[6pt]
\hline
\textnormal{Deck Overhang Width (m)} & 1.5 m \\[6pt]
\hline
\textnormal{Deck Thickness (mm)} & 250 mm \\[6pt]
\hline
\textnormal{Footpath Width (m)} & 1.5 m (IRC 5 Cl. 104.3.6 min: 1.5 m) \\[6pt]
\hline
\textnormal{No. of Traffic Lanes} & 2 (per IRC 5 Cl. 104.3.1) \\[6pt]
\hline
\end{longtable}

\vspace{0.8em}

\begin{longtable}{|L{5.5cm}|p{10.0cm}|}
\caption{\textbf{Components Details}}\label{tab:component-details}\\
\hline
\textbf{Parameter} & \textbf{Value} \\
\hline
\endfirsthead

\multicolumn{2}{c}{\small\textit{\tablename\ \thetable\ -- Continued from previous page}} \\
\hline
\textbf{Parameter} & \textbf{Value} \\
\hline
\endhead

\hline
\multicolumn{2}{r}{\footnotesize\textit{Continued on next page\ldots}} \\
\endfoot

\hline
\endlastfoot

\textnormal{Crash Barrier Type} & Crash Barrier Type \\[6pt]
\hline
\textnormal{Crash Barrier Load (kN/m)} & 1.0 \\[6pt]
\hline

\textnormal{Railing Type} & Railing Type \\[6pt]
\hline
\textnormal{Railing Load (kN/m)} & 0.5 \\[6pt]
\hline
\textnormal{Wearing Course Material} & Bituminous Concrete \\[6pt]
\hline
\textnormal{Wearing Course Thickness (mm)} & 65 mm \\[6pt]
\hline
\end{longtable}


\newpage

\vspace{0.4em}
\noindent
            
\vspace{4pt}
\begin{longtable}{|L{2.2cm}|L{1.8cm}|p{3.8cm}|p{3.8cm}|p{3.8cm}|}
\caption{\textbf{Girder General Information}}\label{tab:girder-gen}\\
\hline
\textbf{Girder} & \textbf{Member ID} & \textbf{Design Mode} & \textbf{Girder Type} & \textbf{Girder Symmetry} \\
\hline
\endfirsthead

\multicolumn{5}{c}{\small\textit{\tablename\ \thetable\ -- Continued from previous page}} \\
\hline
\textbf{Girder} & \textbf{Member ID} & \textbf{Design Mode} & \textbf{Girder Type} & \textbf{Girder Symmetry} \\
\hline
\endhead

\hline
\multicolumn{5}{r}{\footnotesize\textit{Continued on next page\ldots}} \\
\endfoot

\hline
\endlastfoot
G1 & G1M1 & Optimized &  &  \\[8pt]
\hline
G2 & G2M1 & Optimized &  &  \\[8pt]
\hline
G3 & G3M1 & Optimized &  &  \\[8pt]
\hline
G4 & G4M1 & Optimized &  &  \\[8pt]
\hline
\end{longtable}

\vspace{0.6em}

\vspace{4pt}
\begin{longtable}{|L{1.8cm}|L{2.3cm}|L{1.8cm}|p{4.8cm}|p{4.8cm}|}
\caption{\textbf{Girder Section Dimensions}}\label{tab:girder-dims}\\
\hline
\textbf{Girder} & \textbf{Total Depth, D (mm)} & \textbf{Web, $t_w$ (mm)} & \textbf{Top Flange ($b_{tf}, t_{tf}$) mm} & \textbf{Bottom Flange ($b_{bf}, t_{bf}$) mm} \\
\hline
\endfirsthead

\multicolumn{5}{c}{\small\textit{\tablename\ \thetable\ -- Continued from previous page}} \\
\hline
\textbf{Girder} & \textbf{Total Depth, D (mm)} & \textbf{Web, $t_w$ (mm)} & \textbf{Top Flange ($b_{tf}, t_{tf}$) mm} & \textbf{Bottom Flange ($b_{bf}, t_{bf}$) mm} \\
\hline
\endhead

\hline
\multicolumn{5}{r}{\footnotesize\textit{Continued on next page\ldots}} \\
\endfoot

\hline
\endlastfoot
G1 & 1864 mm & 16 mm & 500 mm, 28 mm & 550 mm, 36 mm \\[8pt]
\hline
G2 & 1864 mm & 16 mm & 500 mm, 28 mm & 550 mm, 36 mm \\[8pt]
\hline
G3 & 1864 mm & 16 mm & 500 mm, 28 mm & 550 mm, 36 mm \\[8pt]
\hline
G4 & 1864 mm & 16 mm & 500 mm, 28 mm & 550 mm, 36 mm \\[8pt]
\hline
\end{longtable}

\vspace{0.6em}

\vspace{4pt}
\begin{longtable}{|L{1.4cm}|p{2.2cm}|p{2.2cm}|p{3.0cm}|p{2.4cm}|p{2.2cm}|}
\caption{\textbf{Girder Restraint and Stiffener Details}}\label{tab:girder-stiffeners}\\
\hline
\textbf{Girder} & \textbf{Torsional / Warping Restraint} & \textbf{Web Philosophy} & \textbf{Intermediate Stiffeners} & \textbf{Longitudinal Stiffeners} & \textbf{Bearing Stiffener} \\
\hline
\endfirsthead

\multicolumn{6}{c}{\small\textit{\tablename\ \thetable\ -- Continued from previous page}} \\
\hline
\textbf{Girder} & \textbf{Torsional / Warping Restraint} & \textbf{Web Philosophy} & \textbf{Intermediate Stiffeners} & \textbf{Longitudinal Stiffeners} & \textbf{Bearing Stiffener} \\
\hline
\endhead

\hline
\multicolumn{6}{r}{\footnotesize\textit{Continued on next page\ldots}} \\
\endfoot

\hline
\endlastfoot
G1 & ,  &  & ; Spacing: ; Thickness:  &  & No.: ; Spacing: ; Thickness:  \\[8pt]
\hline
G2 & ,  &  & ; Spacing: ; Thickness:  &  & No.: ; Spacing: ; Thickness:  \\[8pt]
\hline
G3 & ,  &  & ; Spacing: ; Thickness:  &  & No.: ; Spacing: ; Thickness:  \\[8pt]
\hline
G4 & ,  &  & ; Spacing: ; Thickness:  &  & No.: ; Spacing: ; Thickness:  \\[8pt]
\hline
\end{longtable}



\newpage

\vspace{0.4em}
\setlength{\tabcolsep}{4pt}
\setlength\LTleft{0pt}
\setlength\LTright{\fill}

\begin{longtable}{|L{2.2cm}|L{2.2cm}|L{3.0cm}|L{2.5cm}|C{1.8cm}|C{1.8cm}|}
\caption{\textbf{Member Properties: Cross Bracing Details}}\label{tab:cb-details}\\
\hline
\textbf{Location} & \textbf{Member IDs} & \textbf{Type of Bracing} & \textbf{Bracing Section} & \textbf{Spacing (m)} \\
\hline
\endfirsthead

\multicolumn{5}{c}{\small\textit{\tablename\ \thetable\ -- Continued from previous page}} \\
\hline
\textbf{Location} & \textbf{Member IDs} & \textbf{Type of Bracing} & \textbf{Bracing Section} & \textbf{Spacing (m)} \\
\hline
\endhead

\hline
\multicolumn{5}{r}{\footnotesize\textit{Continued on next page\ldots}} \\
\endfoot

\hline
\endlastfoot
 &  &  &  &  \\[6pt]
\hline
 &  &  &  &  \\[6pt]
\hline
 &  &  &  &  \\[6pt]
\hline
\end{longtable}

\vspace{0.4em}
\noindent
\setlength{\tabcolsep}{4pt}
\setlength\LTleft{0pt}
\setlength\LTright{\fill}

\begin{longtable}{|L{2.2cm}|L{2.2cm}|L{3.0cm}|L{2.5cm}|C{1.8cm}|C{1.8cm}|}
\caption{\textbf{Member Properties: End Diaphragm Details}}\label{tab:ed-details}\\
\hline
\textbf{Location} & \textbf{Member IDs} & \textbf{Type of Bracing} & \textbf{Bracing Section} \\
\hline
\endfirsthead

\multicolumn{4}{c}{\small\textit{\tablename\ \thetable\ -- Continued from previous page}} \\
\hline
\textbf{Location} & \textbf{Member IDs} & \textbf{Type of Bracing} & \textbf{Bracing Section} \\
\hline
\endhead

\hline
\multicolumn{4}{r}{\footnotesize\textit{Continued on next page\ldots}} \\
\endfoot

\hline
\endlastfoot
 &  &  &  \\[6pt]
\hline
 &  &  &  \\[6pt]
\hline
 &  &  &  \\[6pt]
\hline
\end{longtable}



\label{subsec:shear-connectors}

\vspace{2.2em}

\vspace{0.4em}
\begin{longtable}{|L{5.5cm}|p{10.0cm}|}
\caption{\textbf{Shear Connector Details}}\label{tab:shear-conn}\\
\hline
\textbf{Parameter} & \textbf{Value} \\
\hline
\endfirsthead

\multicolumn{2}{c}{\small\textit{\tablename\ \thetable\ -- Continued from previous page}} \\
\hline
\textbf{Parameter} & \textbf{Value} \\
\hline
\endhead

\hline
\multicolumn{2}{r}{\footnotesize\textit{Continued on next page\ldots}} \\
\endfoot

\hline
\endlastfoot

\textnormal{Stud Diameter (mm)} & 22 mm \\[6pt]
\hline
\textnormal{Stud Height (mm)} & 150 mm \\[6pt]
\hline
\textnormal{Stud $f_y$ (MPa)} &  \\[6pt]
\hline
\textnormal{Stud $f_u$ (MPa)} &  \\[6pt]
\hline
\textnormal{No. of Studs per Section} & 3 \\[6pt]
\hline
\end{longtable}



\label{subsec:safety-factors}

\vspace{2.2em}

\vspace{0.3em}
\noindent\textit{Note: All values are per IRC 22 Table 1 unless user-modified.}

\vspace{0.4em}
\begin{longtable}{|L{5.5cm}|p{10.0cm}|}
\caption{\textbf{Partial Safety Factors}}\label{tab:safety-factors}\\
\hline
\textbf{Parameter} & \textbf{Value} \\
\hline
\endfirsthead

\multicolumn{2}{c}{\small\textit{\tablename\ \thetable\ -- Continued from previous page}} \\
\hline
\textbf{Parameter} & \textbf{Value} \\
\hline
\endhead

\hline
\multicolumn{2}{r}{\footnotesize\textit{Continued on next page\ldots}} \\
\endfoot

\hline
\endlastfoot

\textnormal{$\gamma_{M0}$ (Yielding / Buckling)} &  \\[6pt]
\hline
\textnormal{$\gamma_{M1}$ (Ultimate Stress)} &  \\[6pt]
\hline
\textnormal{$\gamma_C$ (Concrete, Basic)} &  \\[6pt]
\hline
\textnormal{$\gamma_s$ (Reinforcement)} &  \\[6pt]
\hline
\textnormal{$\gamma_v$ (Shear Connectors)} &  \\[6pt]
\hline
\textnormal{$\gamma_{fft}$ (Fatigue Load)} &  \\[6pt]
\hline
\textnormal{$\gamma_{Mft}$ (Fatigue Strength)} &  \\[6pt]
\hline
\end{longtable}


\chapter{Loads and Load Combinations}
\label{ch:loads-combinations}

This section summarizes all permanent and transient loads applied to the bridge and the governing load combinations evaluated for analysis and design per IRC 6:2017.

\vspace{1em}
\begin{longtable}{|L{6.0cm}|p{9.5cm}|}
\caption{\textbf{Dead Load --- Self Weight}}\label{tab:dl-self-weight}\\
\hline
\textbf{Parameter} & \textbf{Value} \\
\hline
\endfirsthead

\multicolumn{2}{c}{\small\textit{\tablename\ \thetable\ -- Continued from previous page}} \\
\hline
\textbf{Parameter} & \textbf{Value} \\
\hline
\endhead

\hline
\multicolumn{2}{r}{\footnotesize\textit{Continued on next page\ldots}} \\
\endfoot

\hline
\endlastfoot

\textnormal{Steel Self-Weight Applied} & 78.5 kN/m\textsuperscript{3} \\[6pt]
\hline
\textnormal{Concrete Deck Weight} & 25.0 kN/m\textsuperscript{3} \\[6pt]
\hline
\textnormal{Self-Weight Factor} & 1.0 \\[6pt]
\hline
\end{longtable}

\vspace{1em}
\begin{longtable}{|L{6.0cm}|p{9.5cm}|}
\caption{\textbf{Dead Load for Surfacing (DW / SIDL)}}\label{tab:dl-surfacing}\\
\hline
\textbf{Parameter} & \textbf{Value} \\
\hline
\endfirsthead

\multicolumn{2}{c}{\small\textit{\tablename\ \thetable\ -- Continued from previous page}} \\
\hline
\textbf{Parameter} & \textbf{Value} \\
\hline
\endhead

\hline
\multicolumn{2}{r}{\footnotesize\textit{Continued on next page\ldots}} \\
\endfoot

\hline
\endlastfoot

\textnormal{Wearing Course Material} & Bituminous Concrete \\[6pt]
\hline
\textnormal{Wearing Course Thickness} & 65 mm \\[6pt]
\hline
\textnormal{Additional SIDL (Crash Barrier)} & 1.0 kN/m per barrier \\[6pt]
\hline
\textnormal{Railing Load} & 0.5 kN/m\sdstar{} \\[6pt]
\hline
\end{longtable}

\vspace{1em}
\begin{longtable}{|L{6.0cm}|p{9.5cm}|}
\caption{\textbf{Vehicle Live Loads (LL) --- per IRC 6}}\label{tab:vehicle-live-load}\\
\hline
\textbf{Parameter} & \textbf{Value} \\
\hline
\endfirsthead

\multicolumn{2}{c}{\small\textit{\tablename\ \thetable\ -- Continued from previous page}} \\
\hline
\textbf{Parameter} & \textbf{Value} \\
\hline
\endhead

\hline
\multicolumn{2}{r}{\footnotesize\textit{Continued on next page\ldots}} \\
\endfoot

\hline
\endlastfoot

\textnormal{Selected Live Load Vehicle Class} & \textbf{IRC Class A, IRC Class 70R (Wheeled)} \\[6pt]
\hline
\textnormal{Impact Factor ($1.0 + i$)} & Class A:\allowbreak{} 1.207 (fraction = 0.207), Class AA / 70R:\allowbreak{} 1.100 (fraction = 0.100) \\[6pt]
\hline
\textnormal{Braking Load ($F_b$)} & N/A \\[6pt]
\hline
\textnormal{Centrifugal Force ($F_{cf}$)} & Nil / Not Applicable (Straight Alignment per IRC 6 Cl. 212) \\[6pt]
\hline
\textnormal{Number of Traffic Lanes} & 2 (per IRC 5 Cl. 104.3.1) \\[6pt]
\hline
\textnormal{Carriageway Width} & 7.5 m \\[6pt]
\hline
\end{longtable}

\vspace{1em}
\begin{longtable}{|L{6.0cm}|p{9.5cm}|}
\caption{\textbf{Footway Live Loads --- per IRC 6 & IRC 5}}\label{tab:footway-load}\\
\hline
\textbf{Parameter} & \textbf{Value} \\
\hline
\endfirsthead

\multicolumn{2}{c}{\small\textit{\tablename\ \thetable\ -- Continued from previous page}} \\
\hline
\textbf{Parameter} & \textbf{Value} \\
\hline
\endhead

\hline
\multicolumn{2}{r}{\footnotesize\textit{Continued on next page\ldots}} \\
\endfoot

\hline
\endlastfoot

\textnormal{Footway Provision} & Both Sides \\[6pt]
\hline
\textnormal{Footpath Width} & 1.5 m (IRC 5 Cl. 104.3.6 min: 1.5 m) \\[6pt]
\hline
\textnormal{Footway Loading Calculation Mode} & As per IRC 6 \\[6pt]
\hline
\textnormal{Footway Live Load Intensity} & 4.905 kN/m² (IRC 6 Cl. 206.1 code default) \\[6pt]
\hline
\textnormal{Railing Load on Footway} & 0.5 kN/m \\[6pt]
\hline
\end{longtable}

\vspace{1em}
\begin{longtable}{|L{6.0cm}|p{9.5cm}|}
\caption{\textbf{Wind Load (WL) --- per IRC 6}}\label{tab:wind-load}\\
\hline
\textbf{Parameter} & \textbf{Value} \\
\hline
\endfirsthead

\multicolumn{2}{c}{\small\textit{\tablename\ \thetable\ -- Continued from previous page}} \\
\hline
\textbf{Parameter} & \textbf{Value} \\
\hline
\endhead

\hline
\multicolumn{2}{r}{\footnotesize\textit{Continued on next page\ldots}} \\
\endfoot

\hline
\endlastfoot

\textnormal{Basic Wind Speed, $V_b$} & 39.0 m/s [from Project Location] \\[6pt]
\hline
\textnormal{Terrain Type} & Plain Terrain \\[6pt]
\hline
\textnormal{Average Exposed Height, $H$} & 10.0 m \\[6pt]
\hline
\textnormal{Hourly Mean Wind Speed, $V_z$} &  \\[6pt]
\hline
\textnormal{Hourly Wind Pressure, $P_z$} &  \\[6pt]
\hline
\textnormal{Transverse Wind Force} &  \\[6pt]
\hline
\textnormal{Longitudinal Wind Force} &  \\[6pt]
\hline
\textnormal{Vertical Wind Force} &  \\[6pt]
\hline
\end{longtable}

\vspace{1em}
\begin{longtable}{|L{6.0cm}|p{9.5cm}|}
\caption{\textbf{Earthquake Load (EL) --- per IRC 6}}\label{tab:earthquake-load}\\
\hline
\textbf{Parameter} & \textbf{Value} \\
\hline
\endfirsthead

\multicolumn{2}{c}{\small\textit{\tablename\ \thetable\ -- Continued from previous page}} \\
\hline
\textbf{Parameter} & \textbf{Value} \\
\hline
\endhead

\hline
\multicolumn{2}{r}{\footnotesize\textit{Continued on next page\ldots}} \\
\endfoot

\hline
\endlastfoot

\textnormal{Seismic Zone} & Zone IV [from Project Location] \\[6pt]
\hline
\textnormal{Zone Factor, $Z$} & 0.24 \\[6pt]
\hline
\textnormal{Importance Factor, $I$} &  \\[6pt]
\hline
\textnormal{Type of Soil} & Type II – Medium Soil \\[6pt]
\hline
\textnormal{Average Response Acceleration, $S_a/g$} &  \\[6pt]
\hline
\textnormal{Horizontal Seismic Coefficient, $A_h$} &  \\[6pt]
\hline
\textnormal{Vertical Seismic Coefficient, $A_v$} &  \\[6pt]
\hline
\end{longtable}

\vspace{1em}
\begin{longtable}{|L{6.0cm}|p{9.5cm}|}
\caption{\textbf{Temperature Load (TL) --- per IRC 6}}\label{tab:temperature-load}\\
\hline
\textbf{Parameter} & \textbf{Value} \\
\hline
\endfirsthead

\multicolumn{2}{c}{\small\textit{\tablename\ \thetable\ -- Continued from previous page}} \\
\hline
\textbf{Parameter} & \textbf{Value} \\
\hline
\endhead

\hline
\multicolumn{2}{r}{\footnotesize\textit{Continued on next page\ldots}} \\
\endfoot

\hline
\endlastfoot

\textnormal{Maximum Shade Temperature} & 45.0 $^\circ$C \\[6pt]
\hline
\textnormal{Minimum Shade Temperature} & -5.0 $^\circ$C \\[6pt]
\hline
\textnormal{Effective Bridge Temp. Range} &  to  $^\circ$C \\[6pt]
\hline
\textnormal{Temperature Rise / Fall for Design} & + $^\circ$C / \textminus{} $^\circ$C \\[6pt]
\hline
\end{longtable}

\vspace{1em}
\begin{longtable}{|C{3.8cm}|p{11.7cm}|}
\caption{\textbf{Load Combinations (ULS and SLS) --- IRC 6 Table B.1 to B.4}}\label{tab:load-combos}\\
\hline
\textbf{Combination ID} & \textbf{Load Cases and Partial Safety Factors} \\
\hline
\endfirsthead

\multicolumn{2}{c}{\small\textit{\tablename\ \thetable\ -- Continued from previous page}} \\
\hline
\textbf{Combination ID} & \textbf{Load Cases and Partial Safety Factors} \\
\hline
\endhead

\hline
\multicolumn{2}{r}{\footnotesize\textit{Continued on next page\ldots}} \\
\endfoot

\hline
\endlastfoot

ULS-01 & DL(1.35/1.00) + SIDL(1.75/1.00) + LL(1.50) + WL(0.90) + TL(0.90) \\[6pt]
\hline
ULS-02 & DL(1.35/1.00) + SIDL(1.75/1.00) + LL(1.15) + WL(1.50) + TL(0.90) \\[6pt]
\hline
ULS-03 & DL(1.35/1.00) + SIDL(1.75/1.00) + LL(1.15) + WL(0.90) + TL(1.50) \\[6pt]
\hline
ULS-04 & DL(1.00/1.00) + SIDL(1.00/1.00) + LL(0.75) + TL(0.50) + VC(1.00) \\[6pt]
\hline
ULS-05 & DL(1.00/1.00) + SIDL(1.00/1.00) + LL(0.75) + TL(0.50) + BI(1.00) \\[6pt]
\hline
ULS-06 & DL(1.00/1.00) + SIDL(1.00/1.00) + LL(0.75) + TL(0.50) + FB(1.00) \\[6pt]
\hline
ULS-07 & DL(1.35/1.00) + SIDL(1.75/1.00) + LL(0.20) + TL(0.50) + EQ(1.50) \\[6pt]
\hline
ULS-08 & DL(1.35/1.00) + SIDL(1.75/1.00) + LL(0.20) + TL(0.50) + EQ(0.75) \\[6pt]
\hline
SLS-01 & DL(1.00) + SIDL(1.20/1.00) + LL(1.00) + WL(0.60) + TL(0.60) \\[6pt]
\hline
SLS-02 & DL(1.00) + SIDL(1.20/1.00) + LL(0.75) + WL(1.00) + TL(0.60) \\[6pt]
\hline
SLS-03 & DL(1.00) + SIDL(1.20/1.00) + LL(0.75) + WL(0.60) + TL(1.00) \\[6pt]
\hline
SLS-04 & DL(1.00) + SIDL(1.20/1.00) + LL(0.75) + WL(0.50) + TL(0.50) \\[6pt]
\hline
SLS-05 & DL(1.00) + SIDL(1.20/1.00) + LL(0.20) + WL(0.60) + TL(0.50) \\[6pt]
\hline
SLS-06 & DL(1.00) + SIDL(1.20/1.00) + LL(0.20) + WL(0.50) + TL(0.60) \\[6pt]
\hline
SLS-07 & DL(1.00) + SIDL(1.20/1.00) + LL(0.00) + WL(0.00) + TL(0.50) \\[6pt]
\hline
\end{longtable}

\noindent\textit{Note: All IRC 6 load combinations are auto-generated by OsdagBridge. User-defined custom combinations, if any, are appended.}


\chapter{Analysis Results}

A grillage model was used for structural analysis. The deck is idealized as a grid of elastic beam elements --- longitudinal members represent the composite steel girders with effective slab, and transverse members represent the slab or cross frames. This section summarizes the critical output from that analysis.

\vspace{1em}
\begingroup
\footnotesize
\setlength{\tabcolsep}{3pt}
\renewcommand{\arraystretch}{1.25}

\begin{longtable}{|
>{\centering\arraybackslash}p{3.1cm}|
>{\centering\arraybackslash}p{1.7cm}|
>{\centering\arraybackslash}p{1.2cm}|
>{\centering\arraybackslash}p{1.3cm}|
>{\centering\arraybackslash}p{1.7cm}|
>{\centering\arraybackslash}p{1.2cm}|
>{\centering\arraybackslash}p{1.3cm}|
>{\centering\arraybackslash}p{1.6cm}|
>{\centering\arraybackslash}p{1.6cm}|}

\caption{\textbf{Summary of Maximum Demands}}\\
\hline
\multirow{2}{*}{\makecell{\textbf{Load}\\\textbf{Case/}\\\textbf{Comb.}}}
& \multicolumn{3}{c|}{\textbf{Bending Moment}}
& \multicolumn{3}{c|}{\textbf{Shear Force}}
& \multicolumn{2}{c|}{\textbf{Reaction at Supports}}\\
\cline{2-9}

& \makecell{\textbf{Max}\\\textbf{(kNm)}}
& \textbf{Girder}
& \makecell{\textbf{Loc.}\\\textbf{(m)}}
& \makecell{\textbf{Max}\\\textbf{(kN)}}
& \textbf{Girder}
& \makecell{\textbf{Loc.}\\\textbf{(m)}}
& \makecell{\textbf{Left}}
& \makecell{\textbf{Right }}\\
\hline
\endfirsthead

\hline
\multirow{2}{*}{\makecell{\textbf{Load}\\\textbf{Case/}\\\textbf{Comb.}}}
& \multicolumn{3}{c|}{\textbf{Bending Moment}}
& \multicolumn{3}{c|}{\textbf{Shear Force}}
& \multicolumn{2}{c|}{\textbf{Reaction at Supports}}\\
\cline{2-9}

& \makecell{\textbf{Max}\\\textbf{(kNm)}}
& \textbf{Girder}
& \makecell{\textbf{Loc.}\\\textbf{(m)}}
& \makecell{\textbf{Max}\\\textbf{(kN)}}
& \textbf{Girder}
& \makecell{\textbf{Loc.}\\\textbf{(m)}}
& \makecell{\textbf{Left}}
& \makecell{\textbf{Right}}\\
\hline
\endhead

DL & 1850.500 & G2 & 15.000 & 450.200 & G1 & 0.000 & 450.200 & 450.200 \\[6pt]
\hline
SIDL & 620.000 & G2 & 15.000 & 165.400 & G1 & 0.000 & 165.400 & 165.400 \\[6pt]
\hline
LL (Class A) & 1380.000 & G2 & 15.000 & 624.600 & G1 & 0.000 & 624.600 & 624.600 \\[6pt]

\hline
\end{longtable}
\endgroup

\vspace{1em}
\begin{longtable}{|>{\centering\arraybackslash}p{5.2cm}|>{\centering\arraybackslash}p{5.2cm}|>{\centering\arraybackslash}p{5.2cm}|}
\caption{\textbf{Reactions at Supports}}
\hline
\textbf{Load Case} & \textbf{Left Support (kN)} & \textbf{Right Support (kN)} \\[6pt]
\hline
 &  &  \\[6pt]
\hline
 &  &  \\[6pt]
\hline
 &  &  \\[6pt]
\hline
\end{longtable}

\vspace{1em}
\begin{longtable}{|L{7cm}|p{8.5cm}|}
\caption{\textbf{Deflection Summary (Live Load \& Total Load)}}
\hline
\textbf{parameter} & \textbf{value} \\
\hline
\textnormal{Deflection due to Live Load, $\delta_{LL}$} & 18.200 mm \\[6pt]
\hline
\textnormal{Allowable Live Load Deflection ($\Delta_{allow}$)} & L/800 = 37.5 mm \\[6pt]
\hline
\textnormal{Live Load Deflection Check Status} & PASS \\[6pt]
\hline
\textnormal{Deflection due to Total Load, $\delta_{total}$} & 34.500 mm \\[6pt]
\hline
\textnormal{Allowable Total Deflection ($\Delta_{allow}$)} & L/600 = 50.0 mm \\[6pt]
\hline
\textnormal{Total Load Deflection Check Status} & PASS \\[6pt]
\hline
\end{longtable}

\vspace{1em}
\noindent
\noindent\fbox{\parbox{0.97\textwidth}{
\textit{[ PLACEHOLDER: Bending Moment Envelope (Envelope ULS): Max/min BM along span. X-axis: distance from left support (m). Y-axis: Bending Moment (kN-m). ]}
}}
\noindent\fbox{\parbox{0.97\textwidth}{
\textit{[ PLACEHOLDER: Shear Force Envelope (Envelope ULS): Max/min SF along span. X-axis: distance from left support (m). Y-axis: Shear Force (kN). ]}
}}
\noindent\fbox{\parbox{0.97\textwidth}{
\textit{[ PLACEHOLDER: Vertical Deflection D$_y$ (1.0 LL): Maximum deflection along span. Load Case: 1.0 LL, Combination: $D_y$. Nodes shown. Isometric view. ]}
}}


\chapter{Design Checks}

This section presents all structural design checks performed by OsdagBridge. For each member, the demand from the governing load combination, the code-based capacity, and the utilization ratio are tabulated. All checks reference IS 800:2007 and IRC 22:2014 unless stated otherwise.

\section{Plate Girder Design}
\label{sec:plate-girder}

\vspace{1em}
\begin{longtable}{|C{2.5cm}|L{8.0cm}|>{\centering\arraybackslash}p{5.0cm}|}
\caption{\textbf{Girder Section Properties (Final Optimized / User-selected)}}\label{tab:girder-section-properties}\\
\hline
\textbf{Girder} & \textbf{Property} & \textbf{Value} \\
\hline
\endfirsthead

\multicolumn{3}{c}{\small\textit{\tablename\ \thetable\ -- Continued from previous page}} \\
\hline
\textbf{Girder} & \textbf{Property} & \textbf{Value} \\
\hline
\endhead

\hline
\multicolumn{3}{r}{\footnotesize\textit{Continued on next page\ldots}} \\
\endfoot

\hline
\endlastfoot
\multirow{13}{*}{\makecell{G1}} & \textnormal{Depth, D (mm)} & 1864 \\[6pt]
\cline{2-3}
 & \textnormal{Top Flange Width, $b_f$ (mm)} & 500 \\[6pt]
\cline{2-3}
 & \textnormal{Bottom Flange Width, $b_f$ (mm)} & 550 \\[6pt]
\cline{2-3}
 & \textnormal{Top Flange Thickness, $t_f$ (mm)} & 28 \\[6pt]
\cline{2-3}
 & \textnormal{Bottom Flange Thickness, $t_f$ (mm)} & 36 \\[6pt]
\cline{2-3}
 & \textnormal{Web Thickness, $t_w$ (mm)} & 16 \\[6pt]
\cline{2-3}
 & \textnormal{Gross Area, A (cm$^2$)} & 652.8 \\[6pt]
\cline{2-3}
 & \textnormal{Moment of Inertia, $I_z$ (cm$^4$)} & 3824500 \\[6pt]
\cline{2-3}
 & \textnormal{Elastic Section Modulus, $Z_{ez}$ (cm$^3$)} & 41035 \\[6pt]
\cline{2-3}
 & \textnormal{Plastic Section Modulus, $Z_{pz}$ (cm$^3$)} & 46820 \\[6pt]
\cline{2-3}
 & \textnormal{Effective Slab Width, $b_{eff}$ (mm)} & 2150 \\[6pt]
\cline{2-3}
 & \textnormal{Transformed Composite $I_z$ (cm$^4$)} & 6254100 \\[6pt]
\cline{2-3}
 & \textnormal{Depth to Plastic Neutral Axis (mm)} & 220 \\[6pt]
\hline
\multirow{13}{*}{\makecell{G2}} & \textnormal{Depth, D (mm)} & 1864 \\[6pt]
\cline{2-3}
 & \textnormal{Top Flange Width, $b_f$ (mm)} & 500 \\[6pt]
\cline{2-3}
 & \textnormal{Bottom Flange Width, $b_f$ (mm)} & 550 \\[6pt]
\cline{2-3}
 & \textnormal{Top Flange Thickness, $t_f$ (mm)} & 28 \\[6pt]
\cline{2-3}
 & \textnormal{Bottom Flange Thickness, $t_f$ (mm)} & 36 \\[6pt]
\cline{2-3}
 & \textnormal{Web Thickness, $t_w$ (mm)} & 16 \\[6pt]
\cline{2-3}
 & \textnormal{Gross Area, A (cm$^2$)} & 652.8 \\[6pt]
\cline{2-3}
 & \textnormal{Moment of Inertia, $I_z$ (cm$^4$)} & 3824500 \\[6pt]
\cline{2-3}
 & \textnormal{Elastic Section Modulus, $Z_{ez}$ (cm$^3$)} & 41035 \\[6pt]
\cline{2-3}
 & \textnormal{Plastic Section Modulus, $Z_{pz}$ (cm$^3$)} & 46820 \\[6pt]
\cline{2-3}
 & \textnormal{Effective Slab Width, $b_{eff}$ (mm)} & 2150 \\[6pt]
\cline{2-3}
 & \textnormal{Transformed Composite $I_z$ (cm$^4$)} & 6254100 \\[6pt]
\cline{2-3}
 & \textnormal{Depth to Plastic Neutral Axis (mm)} & 220 \\[6pt]
\hline
\multirow{13}{*}{\makecell{G3}} & \textnormal{Depth, D (mm)} & 1864 \\[6pt]
\cline{2-3}
 & \textnormal{Top Flange Width, $b_f$ (mm)} & 500 \\[6pt]
\cline{2-3}
 & \textnormal{Bottom Flange Width, $b_f$ (mm)} & 550 \\[6pt]
\cline{2-3}
 & \textnormal{Top Flange Thickness, $t_f$ (mm)} & 28 \\[6pt]
\cline{2-3}
 & \textnormal{Bottom Flange Thickness, $t_f$ (mm)} & 36 \\[6pt]
\cline{2-3}
 & \textnormal{Web Thickness, $t_w$ (mm)} & 16 \\[6pt]
\cline{2-3}
 & \textnormal{Gross Area, A (cm$^2$)} & 652.8 \\[6pt]
\cline{2-3}
 & \textnormal{Moment of Inertia, $I_z$ (cm$^4$)} & 3824500 \\[6pt]
\cline{2-3}
 & \textnormal{Elastic Section Modulus, $Z_{ez}$ (cm$^3$)} & 41035 \\[6pt]
\cline{2-3}
 & \textnormal{Plastic Section Modulus, $Z_{pz}$ (cm$^3$)} & 46820 \\[6pt]
\cline{2-3}
 & \textnormal{Effective Slab Width, $b_{eff}$ (mm)} & 2150 \\[6pt]
\cline{2-3}
 & \textnormal{Transformed Composite $I_z$ (cm$^4$)} & 6254100 \\[6pt]
\cline{2-3}
 & \textnormal{Depth to Plastic Neutral Axis (mm)} & 220 \\[6pt]
\hline
\multirow{13}{*}{\makecell{G4}} & \textnormal{Depth, D (mm)} & 1864 \\[6pt]
\cline{2-3}
 & \textnormal{Top Flange Width, $b_f$ (mm)} & 500 \\[6pt]
\cline{2-3}
 & \textnormal{Bottom Flange Width, $b_f$ (mm)} & 550 \\[6pt]
\cline{2-3}
 & \textnormal{Top Flange Thickness, $t_f$ (mm)} & 28 \\[6pt]
\cline{2-3}
 & \textnormal{Bottom Flange Thickness, $t_f$ (mm)} & 36 \\[6pt]
\cline{2-3}
 & \textnormal{Web Thickness, $t_w$ (mm)} & 16 \\[6pt]
\cline{2-3}
 & \textnormal{Gross Area, A (cm$^2$)} & 652.8 \\[6pt]
\cline{2-3}
 & \textnormal{Moment of Inertia, $I_z$ (cm$^4$)} & 3824500 \\[6pt]
\cline{2-3}
 & \textnormal{Elastic Section Modulus, $Z_{ez}$ (cm$^3$)} & 41035 \\[6pt]
\cline{2-3}
 & \textnormal{Plastic Section Modulus, $Z_{pz}$ (cm$^3$)} & 46820 \\[6pt]
\cline{2-3}
 & \textnormal{Effective Slab Width, $b_{eff}$ (mm)} & 2150 \\[6pt]
\cline{2-3}
 & \textnormal{Transformed Composite $I_z$ (cm$^4$)} & 6254100 \\[6pt]
\cline{2-3}
 & \textnormal{Depth to Plastic Neutral Axis (mm)} & 220 \\[6pt]
\hline
\end{longtable}

\vspace{1em}
\begin{longtable}{|C{2.5cm}|L{3cm}|C{3.5cm}|C{2.5cm}|>{\centering\arraybackslash}p{4.0cm}|}
\caption{\textbf{Girder Section Classification}}\label{tab:girder-classification}\\
\hline
\textbf{} & \textbf{Element} & \textbf{Slenderness Ratio} & \textbf{Class Limit} & \textbf{Classification} \\
\hline
\endfirsthead

\multicolumn{5}{c}{\small\textit{\tablename\ \thetable\ -- Continued from previous page}} \\
\hline
\textbf{} & \textbf{Element} & \textbf{Slenderness Ratio} & \textbf{Class Limit} & \textbf{Classification} \\
\hline
\endhead

\hline
\multicolumn{5}{r}{\footnotesize\textit{Continued on next page\ldots}} \\
\endfoot

\hline
\endlastfoot
\multirow{3}{*}{\makecell{G1}} & Top Flange & 8.9 & 9.4 & Class 1 (Plastic) \\[6pt]
\cline{2-5}
 & Web & 112.5 & 126.0 & Class 2 (Compact) \\[6pt]
\cline{2-5}
 & Overall Section & --- & --- & Class 2 (Compact) \\[6pt]
\hline
\multirow{3}{*}{\makecell{G2}} & Top Flange & 8.9 & 9.4 & Class 1 (Plastic) \\[6pt]
\cline{2-5}
 & Web & 112.5 & 126.0 & Class 2 (Compact) \\[6pt]
\cline{2-5}
 & Overall Section & --- & --- & Class 2 (Compact) \\[6pt]
\hline
\multirow{3}{*}{\makecell{G3}} & Top Flange & 8.9 & 9.4 & Class 1 (Plastic) \\[6pt]
\cline{2-5}
 & Web & 112.5 & 126.0 & Class 2 (Compact) \\[6pt]
\cline{2-5}
 & Overall Section & --- & --- & Class 2 (Compact) \\[6pt]
\hline
\multirow{3}{*}{\makecell{G4}} & Top Flange & 8.9 & 9.4 & Class 1 (Plastic) \\[6pt]
\cline{2-5}
 & Web & 112.5 & 126.0 & Class 2 (Compact) \\[6pt]
\cline{2-5}
 & Overall Section & --- & --- & Class 2 (Compact) \\[6pt]
\hline
\end{longtable}
\noindent\textit{Note: IS 800:2007 Table 2}

\vspace{1em}
\begin{longtable}{|C{2.5cm}|C{3.5cm}|C{3.5cm}|>{\centering\arraybackslash}p{4.2cm}|C{1.8cm}|}
\caption{\textbf{Moment Capacity Check}}\label{tab:moment-capacity}\\
\hline
\textbf{} & \textbf{Parameter} & \textbf{Formula} & \textbf{Value} & \textbf{Status} \\
\hline
\endfirsthead

\multicolumn{5}{c}{\small\textit{\tablename\ \thetable\ -- Continued from previous page}} \\
\hline
\textbf{} & \textbf{Parameter} & \textbf{Formula} & \textbf{Value} & \textbf{Status} \\
\hline
\endhead

\hline
\multicolumn{5}{r}{\footnotesize\textit{Continued on next page\ldots}} \\
\endfoot

\hline
\endlastfoot
\multirow{3}{*}{\makecell{G1}} & Applied Moment, $M_u$ & Governing LC (ULS) & 3850.5 kN-m & --- \\[6pt]
\cline{2-5}
 & Design Moment Capacity, $M_d$ & IRC 22 Cl. 603.3.1 & 5200.0 kN-m & --- \\[6pt]
\cline{2-5}
 & Utilization Ratio, $M_u / M_d$ & $M_u / M_d$ & 0.74 & PASS \\[6pt]
\hline
\multirow{3}{*}{\makecell{G2}} & Applied Moment, $M_u$ & Governing LC (ULS) & 3850.5 kN-m & --- \\[6pt]
\cline{2-5}
 & Design Moment Capacity, $M_d$ & IRC 22 Cl. 603.3.1 & 5200.0 kN-m & --- \\[6pt]
\cline{2-5}
 & Utilization Ratio, $M_u / M_d$ & $M_u / M_d$ & 0.74 & PASS \\[6pt]
\hline
\multirow{3}{*}{\makecell{G3}} & Applied Moment, $M_u$ & Governing LC (ULS) & 3850.5 kN-m & --- \\[6pt]
\cline{2-5}
 & Design Moment Capacity, $M_d$ & IRC 22 Cl. 603.3.1 & 5200.0 kN-m & --- \\[6pt]
\cline{2-5}
 & Utilization Ratio, $M_u / M_d$ & $M_u / M_d$ & 0.74 & PASS \\[6pt]
\hline
\multirow{3}{*}{\makecell{G4}} & Applied Moment, $M_u$ & Governing LC (ULS) & 3850.5 kN-m & --- \\[6pt]
\cline{2-5}
 & Design Moment Capacity, $M_d$ & IRC 22 Cl. 603.3.1 & 5200.0 kN-m & --- \\[6pt]
\cline{2-5}
 & Utilization Ratio, $M_u / M_d$ & $M_u / M_d$ & 0.74 & PASS \\[6pt]
\hline
\end{longtable}
\noindent\textit{Note: IRC 22 Cl. 603.3.1, IS 800 Cl. 8.2.1}

\vspace{1em}
\begin{longtable}{|C{2.5cm}|C{3.5cm}|C{3.5cm}|>{\centering\arraybackslash}p{4.2cm}|C{1.8cm}|}
\caption{\textbf{Shear Capacity Check}}\label{tab:shear-capacity}\\
\hline
\textbf{} & \textbf{Parameter} & \textbf{Formula} & \textbf{Value} & \textbf{Status} \\
\hline
\endfirsthead

\multicolumn{5}{c}{\small\textit{\tablename\ \thetable\ -- Continued from previous page}} \\
\hline
\textbf{} & \textbf{Parameter} & \textbf{Formula} & \textbf{Value} & \textbf{Status} \\
\hline
\endhead

\hline
\multicolumn{5}{r}{\footnotesize\textit{Continued on next page\ldots}} \\
\endfoot

\hline
\endlastfoot
\multirow{8}{*}{\makecell{G1}} & Applied Shear, $V_u$ & Governing LC (ULS) & 1240.2 kN & --- \\[6pt]
\cline{2-5}
 & Shear Area, $A_v$ & $d_w \times t_w$ & 28800 mm$^2$ & --- \\[6pt]
\cline{2-5}
 & Panel Aspect Ratio, c/d & --- & 1.25 & --- \\[6pt]
\cline{2-5}
 & Shear Buckling Coefficient, $k_v$ & IS 800 Cl. 8.4.2.2 & 5.35 & --- \\[6pt]
\cline{2-5}
 & Web Slenderness, $\lambda_w$ & $\sqrt{f_{yw}/(\sqrt{3}\,\tau_{cr})}$ & 0.88 & --- \\[6pt]
\cline{2-5}
 & Design Shear Stress, $\tau_b$ & IRC 22 Cl. 603.3.3.2 & 118.5 MPa & --- \\[6pt]
\cline{2-5}
 & Shear Buckling Resistance, $V_{cr}$ & $A_v \times \tau_b$ & 1850.0 kN & --- \\[6pt]
\cline{2-5}
 & Utilization Ratio, $V_u / V_d$ & $V_u / V_d$ & 0.67 & PASS \\[6pt]
\hline
\multirow{8}{*}{\makecell{G2}} & Applied Shear, $V_u$ & Governing LC (ULS) & 1240.2 kN & --- \\[6pt]
\cline{2-5}
 & Shear Area, $A_v$ & $d_w \times t_w$ & 28800 mm$^2$ & --- \\[6pt]
\cline{2-5}
 & Panel Aspect Ratio, c/d & --- & 1.25 & --- \\[6pt]
\cline{2-5}
 & Shear Buckling Coefficient, $k_v$ & IS 800 Cl. 8.4.2.2 & 5.35 & --- \\[6pt]
\cline{2-5}
 & Web Slenderness, $\lambda_w$ & $\sqrt{f_{yw}/(\sqrt{3}\,\tau_{cr})}$ & 0.88 & --- \\[6pt]
\cline{2-5}
 & Design Shear Stress, $\tau_b$ & IRC 22 Cl. 603.3.3.2 & 118.5 MPa & --- \\[6pt]
\cline{2-5}
 & Shear Buckling Resistance, $V_{cr}$ & $A_v \times \tau_b$ & 1850.0 kN & --- \\[6pt]
\cline{2-5}
 & Utilization Ratio, $V_u / V_d$ & $V_u / V_d$ & 0.67 & PASS \\[6pt]
\hline
\multirow{8}{*}{\makecell{G3}} & Applied Shear, $V_u$ & Governing LC (ULS) & 1240.2 kN & --- \\[6pt]
\cline{2-5}
 & Shear Area, $A_v$ & $d_w \times t_w$ & 28800 mm$^2$ & --- \\[6pt]
\cline{2-5}
 & Panel Aspect Ratio, c/d & --- & 1.25 & --- \\[6pt]
\cline{2-5}
 & Shear Buckling Coefficient, $k_v$ & IS 800 Cl. 8.4.2.2 & 5.35 & --- \\[6pt]
\cline{2-5}
 & Web Slenderness, $\lambda_w$ & $\sqrt{f_{yw}/(\sqrt{3}\,\tau_{cr})}$ & 0.88 & --- \\[6pt]
\cline{2-5}
 & Design Shear Stress, $\tau_b$ & IRC 22 Cl. 603.3.3.2 & 118.5 MPa & --- \\[6pt]
\cline{2-5}
 & Shear Buckling Resistance, $V_{cr}$ & $A_v \times \tau_b$ & 1850.0 kN & --- \\[6pt]
\cline{2-5}
 & Utilization Ratio, $V_u / V_d$ & $V_u / V_d$ & 0.67 & PASS \\[6pt]
\hline
\multirow{8}{*}{\makecell{G4}} & Applied Shear, $V_u$ & Governing LC (ULS) & 1240.2 kN & --- \\[6pt]
\cline{2-5}
 & Shear Area, $A_v$ & $d_w \times t_w$ & 28800 mm$^2$ & --- \\[6pt]
\cline{2-5}
 & Panel Aspect Ratio, c/d & --- & 1.25 & --- \\[6pt]
\cline{2-5}
 & Shear Buckling Coefficient, $k_v$ & IS 800 Cl. 8.4.2.2 & 5.35 & --- \\[6pt]
\cline{2-5}
 & Web Slenderness, $\lambda_w$ & $\sqrt{f_{yw}/(\sqrt{3}\,\tau_{cr})}$ & 0.88 & --- \\[6pt]
\cline{2-5}
 & Design Shear Stress, $\tau_b$ & IRC 22 Cl. 603.3.3.2 & 118.5 MPa & --- \\[6pt]
\cline{2-5}
 & Shear Buckling Resistance, $V_{cr}$ & $A_v \times \tau_b$ & 1850.0 kN & --- \\[6pt]
\cline{2-5}
 & Utilization Ratio, $V_u / V_d$ & $V_u / V_d$ & 0.67 & PASS \\[6pt]
\hline
\end{longtable}
\noindent\textit{Note: IS 800 Cl. 8.4, IRC 22 Cl. 603.3.3.2}

\vspace{1em}
\begin{longtable}{|C{2.5cm}|C{3.5cm}|C{3.5cm}|>{\centering\arraybackslash}p{4.2cm}|C{1.8cm}|}
\caption{\textbf{Interaction Checks (M-V and M-N)}}\label{tab:interaction-checks}\\
\hline
\textbf{} & \textbf{Check} & \textbf{Condition} & \textbf{Value} & \textbf{Status} \\
\hline
\endfirsthead

\multicolumn{5}{c}{\small\textit{\tablename\ \thetable\ -- Continued from previous page}} \\
\hline
\textbf{} & \textbf{Check} & \textbf{Condition} & \textbf{Value} & \textbf{Status} \\
\hline
\endhead

\hline
\multicolumn{5}{r}{\footnotesize\textit{Continued on next page\ldots}} \\
\endfoot

\hline
\endlastfoot
\multirow{4}{*}{\makecell{G1}} & High Shear Condition? & $V_u > 0.6\,V_d$ & No & --- \\[6pt]
\cline{2-5}
 & Reduced Moment Capacity, $M_{dv}$ & IRC 22 Cl. 603.3.3.3 & 5200.0 kN-m & --- \\[6pt]
\cline{2-5}
 & Interaction Check: $M_u \leq M_{dv}$ & --- & 0.78 & PASS \\[6pt]
\cline{2-5}
 & Interaction Check: $N_u/N_{Rd} + M_u/M_{dv} \leq 1.0$ & N/A & N/A & N/A \\[6pt]
\hline
\multirow{4}{*}{\makecell{G2}} & High Shear Condition? & $V_u > 0.6\,V_d$ & No & --- \\[6pt]
\cline{2-5}
 & Reduced Moment Capacity, $M_{dv}$ & IRC 22 Cl. 603.3.3.3 & 5200.0 kN-m & --- \\[6pt]
\cline{2-5}
 & Interaction Check: $M_u \leq M_{dv}$ & --- & 0.78 & PASS \\[6pt]
\cline{2-5}
 & Interaction Check: $N_u/N_{Rd} + M_u/M_{dv} \leq 1.0$ & N/A & N/A & N/A \\[6pt]
\hline
\multirow{4}{*}{\makecell{G3}} & High Shear Condition? & $V_u > 0.6\,V_d$ & No & --- \\[6pt]
\cline{2-5}
 & Reduced Moment Capacity, $M_{dv}$ & IRC 22 Cl. 603.3.3.3 & 5200.0 kN-m & --- \\[6pt]
\cline{2-5}
 & Interaction Check: $M_u \leq M_{dv}$ & --- & 0.78 & PASS \\[6pt]
\cline{2-5}
 & Interaction Check: $N_u/N_{Rd} + M_u/M_{dv} \leq 1.0$ & N/A & N/A & N/A \\[6pt]
\hline
\multirow{4}{*}{\makecell{G4}} & High Shear Condition? & $V_u > 0.6\,V_d$ & No & --- \\[6pt]
\cline{2-5}
 & Reduced Moment Capacity, $M_{dv}$ & IRC 22 Cl. 603.3.3.3 & 5200.0 kN-m & --- \\[6pt]
\cline{2-5}
 & Interaction Check: $M_u \leq M_{dv}$ & --- & 0.78 & PASS \\[6pt]
\cline{2-5}
 & Interaction Check: $N_u/N_{Rd} + M_u/M_{dv} \leq 1.0$ & N/A & N/A & N/A \\[6pt]
\hline
\end{longtable}
\noindent\textit{Note: IRC 22 Cl. 603.3.3.3}

\vspace{1em}
\begin{longtable}{|C{2.5cm}|C{3.5cm}|C{3.5cm}|>{\centering\arraybackslash}p{4.2cm}|C{1.8cm}|}
\caption{\textbf{Lateral Torsional Buckling Check -- Construction Stage}}\label{tab:ltb-check}\\
\hline
\textbf{} & \textbf{Parameter} & \textbf{Formula} & \textbf{Value} & \textbf{Status} \\
\hline
\endfirsthead

\multicolumn{5}{c}{\small\textit{\tablename\ \thetable\ -- Continued from previous page}} \\
\hline
\textbf{} & \textbf{Parameter} & \textbf{Formula} & \textbf{Value} & \textbf{Status} \\
\hline
\endhead

\hline
\multicolumn{5}{r}{\footnotesize\textit{Continued on next page\ldots}} \\
\endfoot

\hline
\endlastfoot
\multirow{5}{*}{\makecell{G1}} & Elastic Critical Moment, $M_{cr}$ & IRC 22 Cl. 603.3.3.1 &  & --- \\[6pt]
\cline{2-5}
 & Non-dim. Slenderness, $\bar{\lambda}_{LT}$ & $\sqrt{M_p / M_{cr}}$ &  & --- \\[6pt]
\cline{2-5}
 & LTB Reduction Factor, $\chi_{LT}$ & IS 800 Cl. 8.2.2 &  & --- \\[6pt]
\cline{2-5}
 & LTB Resistance, $M_b$ & $\chi_{LT}\,M_p / \gamma_{m0}$ &  & --- \\[6pt]
\cline{2-5}
 & $M_u \leq M_b$ & $M_u / M_b$ & 0.71 & PASS \\[6pt]
\hline
\multirow{5}{*}{\makecell{G2}} & Elastic Critical Moment, $M_{cr}$ & IRC 22 Cl. 603.3.3.1 &  & --- \\[6pt]
\cline{2-5}
 & Non-dim. Slenderness, $\bar{\lambda}_{LT}$ & $\sqrt{M_p / M_{cr}}$ &  & --- \\[6pt]
\cline{2-5}
 & LTB Reduction Factor, $\chi_{LT}$ & IS 800 Cl. 8.2.2 &  & --- \\[6pt]
\cline{2-5}
 & LTB Resistance, $M_b$ & $\chi_{LT}\,M_p / \gamma_{m0}$ &  & --- \\[6pt]
\cline{2-5}
 & $M_u \leq M_b$ & $M_u / M_b$ & 0.71 & PASS \\[6pt]
\hline
\multirow{5}{*}{\makecell{G3}} & Elastic Critical Moment, $M_{cr}$ & IRC 22 Cl. 603.3.3.1 &  & --- \\[6pt]
\cline{2-5}
 & Non-dim. Slenderness, $\bar{\lambda}_{LT}$ & $\sqrt{M_p / M_{cr}}$ &  & --- \\[6pt]
\cline{2-5}
 & LTB Reduction Factor, $\chi_{LT}$ & IS 800 Cl. 8.2.2 &  & --- \\[6pt]
\cline{2-5}
 & LTB Resistance, $M_b$ & $\chi_{LT}\,M_p / \gamma_{m0}$ &  & --- \\[6pt]
\cline{2-5}
 & $M_u \leq M_b$ & $M_u / M_b$ & 0.71 & PASS \\[6pt]
\hline
\multirow{5}{*}{\makecell{G4}} & Elastic Critical Moment, $M_{cr}$ & IRC 22 Cl. 603.3.3.1 &  & --- \\[6pt]
\cline{2-5}
 & Non-dim. Slenderness, $\bar{\lambda}_{LT}$ & $\sqrt{M_p / M_{cr}}$ &  & --- \\[6pt]
\cline{2-5}
 & LTB Reduction Factor, $\chi_{LT}$ & IS 800 Cl. 8.2.2 &  & --- \\[6pt]
\cline{2-5}
 & LTB Resistance, $M_b$ & $\chi_{LT}\,M_p / \gamma_{m0}$ &  & --- \\[6pt]
\cline{2-5}
 & $M_u \leq M_b$ & $M_u / M_b$ & 0.71 & PASS \\[6pt]
\hline
\end{longtable}
\noindent\textit{Note: IRC 22 Cl. 603.3.3.1, IS 800 Cl. 8.2.2}

\vspace{1em}
\begin{longtable}{|C{2.5cm}|L{6.5cm}|>{\arraybackslash}p{6.5cm}|}
\caption{\textbf{Stiffener Design Summary}}\label{tab:stiffener-summary}\\
\hline
\textbf{Girder} & \textbf{Property} & \textbf{Value / Specification} \\
\hline
\endfirsthead

\multicolumn{3}{c}{\small\textit{\tablename\ \thetable\ -- Continued from previous page}} \\
\hline
\textbf{Girder} & \textbf{Property} & \textbf{Value / Specification} \\
\hline
\endhead

\hline
\multicolumn{3}{r}{\footnotesize\textit{Continued on next page\ldots}} \\
\endfoot

\hline
\endlastfoot
\multirow{6}{*}{\makecell{G1}} & \textnormal{Shear Buckling Design Method} &  \\[6pt]
\cline{2-3}
 & \textnormal{Intermediate Stiffener Thickness (mm)} &  \\[6pt]
\cline{2-3}
 & \textnormal{Intermediate Stiffener Spacing (mm)} &  \\[6pt]
\cline{2-3}
 & \textnormal{End Panel Stiffener Thickness (mm)} &  \\[6pt]
\cline{2-3}
 & \textnormal{No. of End Panel Stiffeners} &  \\[6pt]
\cline{2-3}
 & \textnormal{Longitudinal Stiffeners} &  \\[6pt]
\hline
\multirow{6}{*}{\makecell{G2}} & \textnormal{Shear Buckling Design Method} &  \\[6pt]
\cline{2-3}
 & \textnormal{Intermediate Stiffener Thickness (mm)} &  \\[6pt]
\cline{2-3}
 & \textnormal{Intermediate Stiffener Spacing (mm)} &  \\[6pt]
\cline{2-3}
 & \textnormal{End Panel Stiffener Thickness (mm)} &  \\[6pt]
\cline{2-3}
 & \textnormal{No. of End Panel Stiffeners} &  \\[6pt]
\cline{2-3}
 & \textnormal{Longitudinal Stiffeners} &  \\[6pt]
\hline
\multirow{6}{*}{\makecell{G3}} & \textnormal{Shear Buckling Design Method} &  \\[6pt]
\cline{2-3}
 & \textnormal{Intermediate Stiffener Thickness (mm)} &  \\[6pt]
\cline{2-3}
 & \textnormal{Intermediate Stiffener Spacing (mm)} &  \\[6pt]
\cline{2-3}
 & \textnormal{End Panel Stiffener Thickness (mm)} &  \\[6pt]
\cline{2-3}
 & \textnormal{No. of End Panel Stiffeners} &  \\[6pt]
\cline{2-3}
 & \textnormal{Longitudinal Stiffeners} &  \\[6pt]
\hline
\multirow{6}{*}{\makecell{G4}} & \textnormal{Shear Buckling Design Method} &  \\[6pt]
\cline{2-3}
 & \textnormal{Intermediate Stiffener Thickness (mm)} &  \\[6pt]
\cline{2-3}
 & \textnormal{Intermediate Stiffener Spacing (mm)} &  \\[6pt]
\cline{2-3}
 & \textnormal{End Panel Stiffener Thickness (mm)} &  \\[6pt]
\cline{2-3}
 & \textnormal{No. of End Panel Stiffeners} &  \\[6pt]
\cline{2-3}
 & \textnormal{Longitudinal Stiffeners} &  \\[6pt]
\hline
\end{longtable}

\vspace{1em}

\begin{longtable}{|C{2.5cm}|L{3.5cm}|C{3.5cm}|>{\centering\arraybackslash}p{4.2cm}|C{1.8cm}|}
\caption{\textbf{End Panel Stiffener Checks}}\label{tab:end-panel-stiffener}\\
\hline
\textbf{} & \textbf{Check} & \textbf{Required} & \textbf{Provided} & \textbf{Status} \\
\hline
\endfirsthead

\multicolumn{5}{c}{\small\textit{\tablename\ \thetable\ -- Continued from previous page}} \\
\hline
\textbf{} & \textbf{Check} & \textbf{Required} & \textbf{Provided} & \textbf{Status} \\
\hline
\endhead

\hline
\multicolumn{5}{r}{\footnotesize\textit{Continued on next page\ldots}} \\
\endfoot

\hline
\endlastfoot
\multirow{4}{*}{\makecell{G1}} & Web Buckling Resistance & N/A & N/A & N/A \\[6pt]
\cline{2-5}
 & Local Crushing Resistance & N/A & N/A & N/A \\[6pt]
\cline{2-5}
 & Bearing Capacity & N/A & N/A & N/A \\[6pt]
\cline{2-5}
 & Column Buckling Resistance & N/A & N/A & N/A \\[6pt]
\hline
\multirow{4}{*}{\makecell{G2}} & Web Buckling Resistance & N/A & N/A & N/A \\[6pt]
\cline{2-5}
 & Local Crushing Resistance & N/A & N/A & N/A \\[6pt]
\cline{2-5}
 & Bearing Capacity & N/A & N/A & N/A \\[6pt]
\cline{2-5}
 & Column Buckling Resistance & N/A & N/A & N/A \\[6pt]
\hline
\multirow{4}{*}{\makecell{G3}} & Web Buckling Resistance & N/A & N/A & N/A \\[6pt]
\cline{2-5}
 & Local Crushing Resistance & N/A & N/A & N/A \\[6pt]
\cline{2-5}
 & Bearing Capacity & N/A & N/A & N/A \\[6pt]
\cline{2-5}
 & Column Buckling Resistance & N/A & N/A & N/A \\[6pt]
\hline
\multirow{4}{*}{\makecell{G4}} & Web Buckling Resistance & N/A & N/A & N/A \\[6pt]
\cline{2-5}
 & Local Crushing Resistance & N/A & N/A & N/A \\[6pt]
\cline{2-5}
 & Bearing Capacity & N/A & N/A & N/A \\[6pt]
\cline{2-5}
 & Column Buckling Resistance & N/A & N/A & N/A \\[6pt]
\hline
\end{longtable}
\noindent\textit{Note: IS 800 Cl. 8.4.2.2}

\vspace{1em}
\begin{longtable}{|C{2.5cm}|L{3.5cm}|C{3.5cm}|>{\centering\arraybackslash}p{3.5cm}|C{2.5cm}|}
\caption{\textbf{Serviceability -- Deflection Checks}}\label{tab:deflection-checks}\\
\hline
\textbf{} & \textbf{Check} & \textbf{Allowable} & \textbf{Actual} & \textbf{Status} \\
\hline
\endfirsthead

\multicolumn{5}{c}{\small\textit{\tablename\ \thetable\ -- Continued from previous page}} \\
\hline
\textbf{} & \textbf{Check} & \textbf{Allowable} & \textbf{Actual} & \textbf{Status} \\
\hline
\endhead

\hline
\multicolumn{5}{r}{\footnotesize\textit{Continued on next page\ldots}} \\
\endfoot

\hline
\endlastfoot
\multirow{2}{*}{\makecell{G1}} & Live Load Deflection, $\delta_{LL}$ (mm) & L/800 = 37.5 mm & 14.500 & PASS \\[6pt]
\cline{2-5}
 & Total Load Deflection, $\delta_{total}$ (mm) & L/600 = 50.0 mm & 28.200 & PASS \\[6pt]
\hline
\multirow{2}{*}{\makecell{G2}} & Live Load Deflection, $\delta_{LL}$ (mm) & L/800 = 37.5 mm & 18.200 & PASS \\[6pt]
\cline{2-5}
 & Total Load Deflection, $\delta_{total}$ (mm) & L/600 = 50.0 mm & 34.500 & PASS \\[6pt]
\hline
\multirow{2}{*}{\makecell{G3}} & Live Load Deflection, $\delta_{LL}$ (mm) & L/800 = 37.5 mm &  & --- \\[6pt]
\cline{2-5}
 & Total Load Deflection, $\delta_{total}$ (mm) & L/600 = 50.0 mm &  & --- \\[6pt]
\hline
\multirow{2}{*}{\makecell{G4}} & Live Load Deflection, $\delta_{LL}$ (mm) & L/800 = 37.5 mm &  & --- \\[6pt]
\cline{2-5}
 & Total Load Deflection, $\delta_{total}$ (mm) & L/600 = 50.0 mm &  & --- \\[6pt]
\hline
\end{longtable}
\noindent\textit{Note: IRC 22 Cl. 604.3.2}

\vspace{1em}
\begin{longtable}{|C{2.5cm}|L{3.5cm}|C{3.5cm}|>{\centering\arraybackslash}p{3.5cm}|C{2.5cm}|}
\caption{\textbf{Serviceability -- Maximum Stress Limitation}}\label{tab:stress-limitation}\\
\hline
\textbf{} & \textbf{Element} & \textbf{Allowable Stress} & \textbf{Actual Stress} & \textbf{Status} \\
\hline
\endfirsthead

\multicolumn{5}{c}{\small\textit{\tablename\ \thetable\ -- Continued from previous page}} \\
\hline
\textbf{} & \textbf{Element} & \textbf{Allowable Stress} & \textbf{Actual Stress} & \textbf{Status} \\
\hline
\endhead

\hline
\multicolumn{5}{r}{\footnotesize\textit{Continued on next page\ldots}} \\
\endfoot

\hline
\endlastfoot
\makecell{G1} & Structural Steel ($0.9\,f_y$) &  &  & --- \\[6pt]
\hline
\makecell{G2} & Structural Steel ($0.9\,f_y$) &  &  & --- \\[6pt]
\hline
\makecell{G3} & Structural Steel ($0.9\,f_y$) &  &  & --- \\[6pt]
\hline
\makecell{G4} & Structural Steel ($0.9\,f_y$) &  &  & --- \\[6pt]
\hline
\end{longtable}

\vspace{1em}
\begin{longtable}{|C{2.5cm}|C{3.5cm}|C{3.5cm}|>{\centering\arraybackslash}p{3.5cm}|C{2.5cm}|}
\caption{\textbf{Serviceability -- Fatigue Assessment}}\label{tab:fatigue-assessment}\\
\hline
\textbf{} & \textbf{Stress Range, $\Delta\sigma$ (MPa)} & \textbf{Fatigue Limit, $f_{fd}$ (MPa)} & \textbf{Utilization Ratio} & \textbf{Status} \\
\hline
\endfirsthead

\multicolumn{5}{c}{\small\textit{\tablename\ \thetable\ -- Continued from previous page}} \\
\hline
\textbf{} & \textbf{Stress Range, $\Delta\sigma$ (MPa)} & \textbf{Fatigue Limit, $f_{fd}$ (MPa)} & \textbf{Utilization Ratio} & \textbf{Status} \\
\hline
\endhead

\hline
\multicolumn{5}{r}{\footnotesize\textit{Continued on next page\ldots}} \\
\endfoot

\hline
\endlastfoot
\makecell{G1} &  &  &  & --- \\[6pt]
\hline
\makecell{G2} &  &  &  & --- \\[6pt]
\hline
\makecell{G3} &  &  &  & --- \\[6pt]
\hline
\makecell{G4} &  &  &  & --- \\[6pt]
\hline
\end{longtable}
\noindent\textit{Note: IRC 22 Cl. 605 --- governing of normal and shear fatigue (worst by DCR). Capacity reduction factor $\mu_r$ applied where plate thickness > 25 mm.}

\vspace{1em}
\vspace{0.4em}
\begin{longtable}{|C{1.6cm}|>{\centering\arraybackslash}p{3.6cm}|C{2.4cm}|C{2.0cm}|C{2.1cm}|C{1.7cm}|C{1.5cm}|}
\caption{\textbf{Girder Design Summary (DCR / Utilization Ratio)}}\label{tab:girder-dcr-summary}\\
\hline
\textbf{Girder} & \textbf{Controlling LC / Combination} & \textbf{Controlling Check} & \textbf{Demand} & \textbf{Capacity} & \textbf{UR} & \textbf{Status} \\
\hline
\endfirsthead

\multicolumn{7}{c}{\small\textit{\tablename\ \thetable\ -- Continued from previous page}} \\
\hline
\textbf{Girder} & \textbf{Controlling LC / Combination} & \textbf{Controlling Check} & \textbf{Demand} & \textbf{Capacity} & \textbf{UR} & \textbf{Status} \\
\hline
\endhead

\hline
\multicolumn{7}{r}{\footnotesize\textit{Continued on next page\ldots}} \\
\endfoot

\hline
\endlastfoot
G1 &  & Flexure &  &  & 0.740 & PASS \\[6pt]
\hline
G2 &  & Flexure &  &  & 0.740 & PASS \\[6pt]
\hline
G3 &  & Flexure &  &  & 0.740 & PASS \\[6pt]
\hline
G4 &  & Flexure &  &  & 0.740 & PASS \\[6pt]
\hline
\end{longtable}
\noindent\textit{Note: UR = Demand / Capacity. A value $\leq 1.0$ indicates a passing check. The controlling check is the criterion with the highest UR for each girder, with the real load case/combination that drives it.}

\vspace{1em}

\begin{longtable}{|L{3.6cm}|C{5.6cm}|>{\centering\arraybackslash}p{2.6cm}|L{3.0cm}|}
\caption{\textbf{Shear Connector Capacity}}\label{tab:shear-connector-capacity}\\
\hline
\textbf{Parameter} & \textbf{Formula} & \textbf{Value} & \textbf{Reference} \\
\hline
\endfirsthead

\multicolumn{4}{c}{\small\textit{\tablename\ \thetable\ -- Continued from previous page}} \\
\hline
\textbf{Parameter} & \textbf{Formula} & \textbf{Value} & \textbf{Reference} \\
\hline
\endhead

\hline
\multicolumn{4}{r}{\footnotesize\textit{Continued on next page\ldots}} \\
\endfoot

\hline
\endlastfoot
Design Resistance, $Q_u$ & \footnotesize\makecell{$Q_u=\min(Q_{u,s},\,Q_{u,c})$\\[3pt]$Q_{u,s}=\dfrac{0.8\,f_u\,(\pi d^2/4)}{\gamma_v}$\\[3pt]$Q_{u,c}=\dfrac{0.29\,\alpha\,d^2\sqrt{f_{ck}\,E_{cm}}}{\gamma_v}$} &  & IRC 22 Cl. 606.3.1 (Eq. 6.1) \\[6pt]
\hline
Fatigue Shear Resistance, $Q_r$ & IRC 22 Table 8 ($\phi d$, $N_{sc}$) &  & IRC 22 Cl. 606.3.2 (Table 8) \\[6pt]
\hline
\end{longtable}

\vspace{1em}

\setlength\LTleft{0pt}
\setlength\LTright{\fill}

\begin{longtable}{|L{3.2cm}|>{\centering\arraybackslash}p{4.3cm}|>{\centering\arraybackslash}p{4.3cm}|C{2.0cm}|}
\caption{\textbf{Shear Connector Spacing}}\label{tab:shear-connector-spacing}\\
\hline
\textbf{Criterion} & \textbf{Governing Spacing} & \textbf{Actual Spacing Provided} & \textbf{Status} \\
\hline
\endfirsthead

\multicolumn{4}{c}{\small\textit{\tablename\ \thetable\ -- Continued from previous page}} \\
\hline
\textbf{Criterion} & \textbf{Governing Spacing} & \textbf{Actual Spacing Provided} & \textbf{Status} \\
\hline
\endhead

\hline
\multicolumn{4}{r}{\footnotesize\textit{Continued on next page\ldots}} \\
\endfoot

\hline
\endlastfoot
ULS Shear (SL1) &  &  & --- \\[6pt]
\hline
Full Composite (SL2) &  &  & --- \\[6pt]
\hline
SLS Fatigue (SR) &  &  & --- \\[6pt]
\hline
Max Spacing Limit (IRC 22) &  &  & --- \\[6pt]
\hline
\end{longtable}
\noindent\textit{Note: IRC 22 Cl. 606.4, 606.9. Governing spacing $= \min(S_{L1}, S_{L2}, S_R)$.}

\vspace{1em}
\begin{longtable}{|L{5.3cm}|>{\arraybackslash}p{7.2cm}|C{2.0cm}|}
\caption{\textbf{Transverse Shear and Detailing Checks}}\label{tab:transverse-shear-checks}\\
\hline
\textbf{Check} & \textbf{Value} & \textbf{Status} \\
\hline
\endfirsthead

\multicolumn{3}{c}{\small\textit{\tablename\ \thetable\ -- Continued from previous page}} \\
\hline
\textbf{Check} & \textbf{Value} & \textbf{Status} \\
\hline
\endhead

\hline
\multicolumn{3}{r}{\footnotesize\textit{Continued on next page\ldots}} \\
\endfoot

\hline
\endlastfoot
\textnormal{Longitudinal Shear per unit length, $V_L$} &  & --- \\[6pt]
\hline
\textnormal{Transverse Shear Capacity of Slab, $V_{Rd}$} &  & --- \\[6pt]
\hline
\textnormal{Transverse Shear Check} & $V_L/V_{Rd}$ = --- & --- \\[6pt]
\hline
\textnormal{Min. Transverse Reinforcement, $A_{st,min}$} & Required , Provided 0.00 cm$^2$/m & --- \\[6pt]
\hline
\textnormal{Stud Diameter $\leq 2\,t_f$} & $d$ =  $\leq 2t_f$ =  & --- \\[6pt]
\hline
\textnormal{Stud Edge Distance} & Provided  (req. $\geq$ ) & --- \\[6pt]
\hline
\end{longtable}
\noindent\textit{Note: IRC 22 Cl. 606.6, 606.10.}

% ===========================
\section{Deck Slab Design}
\label{sec:deck-design}
% ===========================

The reinforced concrete deck slab is designed per IRC~112:2011 (flexure, shear, crack width) and IRC~22:2014 (composite construction). Wheel loads are distributed using Pigeaud's method. The deck is checked for flexure in the transverse and longitudinal directions, punching shear, one-way (beam) shear, crack width, and reinforcement detailing.

\vspace{1em}
\begin{longtable}{|L{5.5cm}|p{10.0cm}|}
\caption{\textbf{Deck Slab --- Loading and Geometry}}\label{tab:deck-loading-geometry}\\
\hline
\textbf{Parameter} & \textbf{Value / Specification} \\
\hline
\endfirsthead

\multicolumn{2}{c}{\small\textit{\tablename\ \thetable\ -- Continued from previous page}} \\
\hline
\textbf{Parameter} & \textbf{Value / Specification} \\
\hline
\endhead

\hline
\multicolumn{2}{r}{\footnotesize\textit{Continued on next page\ldots}} \\
\endfoot

\hline
\endlastfoot

\textnormal{Effective Span of Deck Slab, $l_{eff}$} & \placeholder{---} mm (girder spacing, c/c) \\[6pt]
\hline
\textnormal{Deck Thickness, $t_s$} & 250 mm \\[6pt]
\hline
\textnormal{Clear Cover (IRC 112 Cl. 15.2)} & Top  / Bottom  mm \\[6pt]
\hline
\textnormal{Concrete Grade (IRC 112 Cl. 6.4)} & M35 ($f_{ck}$ =  MPa, $f_{ctm}$ =  MPa) \\[6pt]
\hline
\textnormal{Reinforcement Grade (IRC 112 Cl. 6.2)} &  ($f_y$ = \placeholder{---} MPa) \\[6pt]
\hline
\textnormal{Dead Load per Unit Area, $w_{DL}$} & \placeholder{---} kN/m² (slab self-weight) \\[6pt]
\hline
\textnormal{IRC 6 Wheel Load (Class A / 70R)} & \placeholder{---} kN \\[6pt]
\hline
\textnormal{Tyre Contact Width (IRC 6 Annex~A)} & \placeholder{---} mm (transverse) \\[6pt]
\hline
\textnormal{Impact Factor (IRC 6 Cl. 208.2)} & \placeholder{---} \\[6pt]
\hline
\textnormal{Governing Live Load Case} & \placeholder{---} \\[6pt]
\hline
\end{longtable}

\vspace{1em}
\begin{longtable}{|C{3.0cm}|C{3.5cm}|C{3.0cm}|>{\centering\arraybackslash}p{4.2cm}|C{1.8cm}|}
\caption{\textbf{Deck Slab --- Flexure Check: Interior Panel (Pigeaud's Method)}}\label{tab:deck-flexure-interior}\\
\hline
\textbf{Location} & \textbf{Parameter} & \textbf{Formula / Reference} & \textbf{Value} & \textbf{Status} \\
\hline
\endfirsthead

\multicolumn{5}{c}{\small\textit{\tablename\ \thetable\ -- Continued from previous page}} \\
\hline
\textbf{Location} & \textbf{Parameter} & \textbf{Formula / Reference} & \textbf{Value} & \textbf{Status} \\
\hline
\endhead

\hline
\multicolumn{5}{r}{\footnotesize\textit{Continued on next page\ldots}} \\
\endfoot

\hline
\endlastfoot

\multirow{5}{*}{\makecell{At Midspan\\(Sagging)}} & Transverse BM (DL), $M_{T,DL}$ & $w_{DL}\,l_{eff}^2/10$ & \placeholder{---} kN-m/m & --- \\[6pt]
\cline{2-5}
 & Transverse BM (LL), $M_{T,LL}$ & Effective width (IRC 112 B3.1) & \placeholder{---} kN-m/m & --- \\[6pt]
\cline{2-5}
 & Total Design BM, $M_{u,sag}$ & \placeholder{---} DL + \placeholder{---} LL & 145.50 kN-m/m & --- \\[6pt]
\cline{2-5}
 & Effective depth, $d$ & $t_s - c_{nom} - \phi/2$ & \placeholder{---} mm & --- \\[6pt]
\cline{2-5}
 & Moment Capacity, $M_{Rd}$ & IRC 112 Cl. 12.2 & 185.00 kN-m/m & PASS \\[6pt]
\hline
\multirow{3}{*}{\makecell{At Support\\(Hogging)}} & Total Design BM, $M_{u,hog}$ & \placeholder{---} DL + \placeholder{---} LL (at support) & 120.00 kN-m/m & --- \\[6pt]
\cline{2-5}
 & Required Top Steel, $A_{st,top}$ & $M_u / (0.87\,f_y\,d)$ & \placeholder{---} mm²/m & --- \\[6pt]
\cline{2-5}
 & Moment Capacity, $M_{Rd}$ & IRC 112 Cl. 12.2 & 160.00 kN-m/m & PASS \\[6pt]
\hline
\end{longtable}
\noindent\textit{Note: IRC 112 Cl. 12.2. Distribution (longitudinal) reinforcement designed for 20\% of main steel moment (IRC 21 Cl. 305.18).}

\vspace{1em}
\begin{longtable}{|L{5.5cm}|C{3.5cm}|>{\centering\arraybackslash}p{4.5cm}|C{2cm}|}
\caption{\textbf{Deck Slab --- Cantilever Overhang Flexure Check}}\label{tab:deck-overhang-flexure}\\
\hline
\textbf{Parameter} & \textbf{Formula} & \textbf{Value} & \textbf{Status} \\
\hline
\endfirsthead

\multicolumn{4}{c}{\small\textit{\tablename\ \thetable\ -- Continued from previous page}} \\
\hline
\textbf{Parameter} & \textbf{Formula} & \textbf{Value} & \textbf{Status} \\
\hline
\endhead

\hline
\multicolumn{4}{r}{\footnotesize\textit{Continued on next page\ldots}} \\
\endfoot

\hline
\endlastfoot

Overhang Length, $l_{oh}$ & --- & 1.5 m & --- \\[6pt]
\hline
Crash Barrier Load Moment & IRC 6 Cl. 206.4 & N/A & --- \\[6pt]
\hline
Dead Load Moment & $w_{DL}\,l_{oh}^2/2$ + railing & N/A & --- \\[6pt]
\hline
Live Load Moment (eccentric wheel) & Wheel load $\times$ arm & N/A & --- \\[6pt]
\hline
Total Hogging Moment, $M_{u,oh}$ & \placeholder{---} DL + \placeholder{---} (LL + CB) & N/A & --- \\[6pt]
\hline
Moment Capacity (top steel), $M_{Rd,oh}$ & IRC 112 Cl. 12.2 & N/A & N/A \\[6pt]
\hline
\end{longtable}
\noindent\textit{Note: IRC 6 Cl. 206.4 crash barrier loads applied at kerb face; IRC 112 Cl. 12.2 flexure.}

\vspace{1em}
\begin{longtable}{|L{5.5cm}|C{3.5cm}|>{\centering\arraybackslash}p{4.5cm}|C{2cm}|}
\caption{\textbf{Deck Slab --- Punching Shear Check (IRC~112 Cl.~10.4.6)}}\label{tab:deck-punching-shear}\\
\hline
\textbf{Parameter} & \textbf{Formula / Reference} & \textbf{Value} & \textbf{Status} \\
\hline
\endfirsthead

\multicolumn{4}{c}{\small\textit{\tablename\ \thetable\ -- Continued from previous page}} \\
\hline
\textbf{Parameter} & \textbf{Formula / Reference} & \textbf{Value} & \textbf{Status} \\
\hline
\endhead

\hline
\multicolumn{4}{r}{\footnotesize\textit{Continued on next page\ldots}} \\
\endfoot

\hline
\endlastfoot

Design Wheel Load (ULS), $V_{Ed}$ & $\gamma_Q\,(1+IF)\,P_w$ & \placeholder{---} kN & --- \\[6pt]
\hline
Tyre Contact Area & $a \times b$ (IRC 6 Annex~A) & \placeholder{---} $\times$ \placeholder{---} mm & --- \\[6pt]
\hline
Loaded Area at mid-depth, $b_0$ & $c_1 \times c_2$ (incl.\ WC dispersion) & \placeholder{---} $\times$ \placeholder{---} mm & --- \\[6pt]
\hline
Control Perimeter, $u_1$ & $2(c_1+c_2) + 4\pi d$ & \placeholder{---} mm & --- \\[6pt]
\hline
Punching Shear Stress, $v_{Ed}$ & $V_{Ed} / (u_1\,d)$ & 245.000 MPa & --- \\[6pt]
\hline
Punching Resistance, $v_{Rd,c}$ & IRC 112 Eq.\ 10.1 & 420.000 MPa & --- \\[6pt]
\hline
Punching Shear Check & $v_{Ed} \leq v_{Rd,c}$ & 0.58 & FAIL \\[6pt]
\hline
\end{longtable}
\noindent\textit{Note: Punching shear reinforcement not typically required for deck slabs with $d \geq 200$ mm and adequate longitudinal reinforcement.}

\vspace{1em}
\clearpage
\begin{longtable}{|L{7cm}|>{\arraybackslash}p{8.5cm}|}
\caption{\textbf{Crack Width Check (Deck Slab)}}\label{tab:deck-crack-width}\\
\hline
\textbf{Parameter} & \textbf{Value / Reference} \\
\hline
\endfirsthead

\multicolumn{2}{c}{\small\textit{\tablename\ \thetable\ -- Continued from previous page}} \\
\hline
\textbf{Parameter} & \textbf{Value / Reference} \\
\hline
\endhead

\hline
\multicolumn{2}{r}{\footnotesize\textit{Continued on next page\ldots}} \\
\endfoot

\hline
\endlastfoot

\textnormal{Min. Reinforcement for Crack Control, $A_{s,min}$} & \placeholder{---} mm²/m [IRC 112 Cl. 16.5.1] \\[6pt]
\hline
\textnormal{Provided Reinforcement (bottom)} & $\phi$\placeholder{---} @ \placeholder{---} mm c/c (\placeholder{---} mm²/m) \\[6pt]
\hline
\textnormal{Max. Permissible Crack Width} & 0.25 mm \\[6pt]
\hline
\textnormal{Calculated Crack Width, $w_k$ (governing)} & 0.1800 mm \\[6pt]
\hline
\textnormal{Crack Width Check} & PASS \\[6pt]
\hline
\end{longtable}

\vspace{1em}
\begin{longtable}{|L{5.5cm}|C{3.5cm}|>{\centering\arraybackslash}p{4.5cm}|C{2cm}|}
\caption{\textbf{One-Way (Beam) Shear Check (Deck Slab)}}\label{tab:deck-beam-shear}\\
\hline
\textbf{Parameter} & \textbf{Formula / Reference} & \textbf{Value} & \textbf{Status} \\
\hline
\endfirsthead

\multicolumn{4}{c}{\small\textit{\tablename\ \thetable\ -- Continued from previous page}} \\
\hline
\textbf{Parameter} & \textbf{Formula / Reference} & \textbf{Value} & \textbf{Status} \\
\hline
\endhead

\hline
\multicolumn{4}{r}{\footnotesize\textit{Continued on next page\ldots}} \\
\endfoot

\hline
\endlastfoot

Design Shear per unit width, $V_{Ed}$ & $\gamma_{DL} V_{DL} + \gamma_{LL}(1{+}IF)V_{LL}$ & 165.00 kN/m & --- \\[6pt]
\hline
Effective depth, $d$ & $t_s - c_{nom} - \phi/2$ & \placeholder{---} mm & --- \\[6pt]
\hline
Size factor, $k$ & $1 + \sqrt{200/d} \leq 2.0$ & \placeholder{---} & --- \\[6pt]
\hline
Long.\ reinforcement ratio, $\rho_l$ & $A_{sl}/(b_w\,d) \leq 0.02$ & \placeholder{---} & --- \\[6pt]
\hline
Shear resistance (no stirrups), $V_{Rd,c}$ & $v_{Rd,c}\,b_w\,d$ (Cl.\ 10.3.2) & 240.00 kN/m & --- \\[6pt]
\hline
One-Way Shear Check & $V_{Ed} \leq V_{Rd,c}$ & 0.69 & FAIL \\[6pt]
\hline
\end{longtable}
\noindent\textit{Note: IRC 112 Cl. 10.3.2. Shear reinforcement not provided in deck slabs; capacity relies on concrete and main reinforcement.}

\vspace{1em}
\begin{longtable}{|L{5.5cm}|>{\centering\arraybackslash}p{4.1cm}|>{\centering\arraybackslash}p{4.1cm}|C{1.8cm}|}
\caption{\textbf{Reinforcement Detailing Summary (Deck Slab)}}\label{tab:deck-rebar-summary}\\
\hline
\textbf{Parameter} & \textbf{Required / Limit} & \textbf{Provided} & \textbf{Status} \\
\hline
\endfirsthead

\multicolumn{4}{c}{\small\textit{\tablename\ \thetable\ -- Continued from previous page}} \\
\hline
\textbf{Parameter} & \textbf{Required / Limit} & \textbf{Provided} & \textbf{Status} \\
\hline
\endhead

\hline
\multicolumn{4}{r}{\footnotesize\textit{Continued on next page\ldots}} \\
\endfoot

\hline
\endlastfoot

\multicolumn{4}{|l|}{\textbf{Main Reinforcement --- Bottom (Transverse)}} \\[6pt]
\hline
Required Area, $A_{st,req}$ (mm²/m) & \placeholder{---} mm²/m & \placeholder{---} mm²/m & PASS \\[6pt]
\hline
Bar Diameter $\times$ Spacing & $\phi \geq 10$ mm (IRC 112) & $\phi$\placeholder{---} @ \placeholder{---} mm c/c & --- \\[6pt]
\hline
Min.\ Reinforcement $A_{s,min}$ (IRC 112 Cl. 16.3.1) & \placeholder{---} mm²/m & \placeholder{---} mm²/m & PASS \\[6pt]
\hline
Max.\ Bar Spacing (IRC 112 Cl. 16.3.2) & \placeholder{---} mm & \placeholder{---} mm & FAIL \\[6pt]
\hline
\multicolumn{4}{|l|}{\textbf{Distribution Reinforcement --- Longitudinal}} \\[6pt]
\hline
Required Area, $A_{st,dist}$ (mm²/m) & $\geq 20\%$ of main steel & \placeholder{---} mm²/m & PASS \\[6pt]
\hline
\multicolumn{4}{|l|}{\textbf{Top Reinforcement (Support / Cantilever Overhang)}} \\[6pt]
\hline
Required Area, $A_{st,top}$ (mm²/m) & \placeholder{---} mm²/m & \placeholder{---} mm²/m & PASS \\[6pt]
\hline
\multicolumn{4}{|l|}{\textbf{Cover and Detailing}} \\[6pt]
\hline
Clear Cover (IRC 112 Cl. 15.2) & $\geq$ \placeholder{---} mm (Table 14.2) & Top  / Bottom  mm & FAIL \\[6pt]
\hline
\end{longtable}
\noindent\textit{Note: IRC 112 Cl. 16.3, IS 456 Cl. 26.5. All reinforcement provisions satisfy strength and detailing requirements.}

% ===========================
\section{Cross Bracing Design}
\label{sec:cross-bracing}
% ===========================

Cross bracing between adjacent plate girders provides lateral stability during construction, resists transverse loads (wind, seismic, braking) in service, and prevents lateral torsional buckling of the girders. Members are designed per IS~800:2007 Cl.~7 (compression) and Cl.~6 (tension). Forces are derived from the grillage model under the governing load combination  (DL + LL + WL).

\vspace{1em}

\vspace{0.4em}
\noindent
\setlength{\tabcolsep}{4pt}
\setlength\LTleft{0pt}
\setlength\LTright{\fill}

\begin{longtable}{|C{2.0cm}|L{2.0cm}|L{2.2cm}|C{2.5cm}|C{2.0cm}|C{2.0cm}|}
\caption{\textbf{Cross Bracing --- Connection and Section Properties}}\label{tab:cb-section-props}\\
\hline
\textbf{Panel} & \textbf{Member} & \textbf{Connection} & \textbf{Section} & \textbf{$A_g$ (mm$^2$)} & \textbf{$r_{min}$ (mm)} \\
\hline
\endfirsthead

\multicolumn{6}{c}{\small\textit{\tablename\ \thetable\ -- Continued from previous page}} \\
\hline
\textbf{Panel} & \textbf{Member} & \textbf{Connection} & \textbf{Section} & \textbf{$A_g$ (mm$^2$)} & \textbf{$r_{min}$ (mm)} \\
\hline
\endhead

\hline
\multicolumn{6}{r}{\footnotesize\textit{Continued on next page\ldots}} \\
\endfoot

\hline
\endlastfoot
\multirow{2}{*}{\makecell{G1-G2}} & Diagonal & Bolted &  & 1550.0 &  \\[6pt]\cline{2-6}
\multirow{2}{*}{\makecell{G1-G2}} & Top / Bottom chord & Bolted &  &  &  \\[6pt]\cline{2-6}
\hline
\multirow{2}{*}{\makecell{G2-G3}} & Diagonal &  &  &  &  \\[6pt]\cline{2-6}
\multirow{2}{*}{\makecell{G2-G3}} & Top / Bottom chord &  &  &  &  \\[6pt]\cline{2-6}
\hline
\multirow{2}{*}{\makecell{G3-G4}} & Diagonal &  &  &  &  \\[6pt]\cline{2-6}
\multirow{2}{*}{\makecell{G3-G4}} & Top / Bottom chord &  &  &  &  \\[6pt]\cline{2-6}
\hline
\end{longtable}
\noindent\textit{Note: $A_g$ = gross cross-sectional area; $r_{min}$ = minimum radius of gyration.}

\vspace{1em}
\begin{longtable}{|C{2.2cm}|L{2.2cm}|L{2.5cm}|C{2.5cm}|C{2.5cm}|>{\centering\arraybackslash}p{3.6cm}|}
\caption{\textbf{Cross Bracing --- Slenderness Ratio Check (IS~800 Cl.~3.8)}}\label{tab:cb-slenderness}\\
\hline
\textbf{Panel} & \textbf{Member} & \textbf{Nature} & \textbf{Eff.\ Length $KL$ (mm)} & \textbf{$KL/r$} & \textbf{Limit / Status} \\
\hline
\endfirsthead

\multicolumn{6}{c}{\small\textit{\tablename\ \thetable\ -- Continued from previous page}} \\
\hline
\textbf{Panel} & \textbf{Member} & \textbf{Nature} & \textbf{Eff.\ Length $KL$ (mm)} & \textbf{$KL/r$} & \textbf{Limit / Status} \\
\hline
\endhead

\hline
\multicolumn{6}{r}{\footnotesize\textit{Continued on next page\ldots}} \\
\endfoot

\hline
\endlastfoot
\multirow{3}{*}{\makecell{G1-G2}} & Diagonal & C & 3120 &  & 250 ---  \\[6pt]
\cline{2-6}
 & Top chord & C & 2500 &  & 250 ---  \\[6pt]
\cline{2-6}
 & Bottom chord & T & 2500 &  & 400 ---  \\[6pt]
\hline
\multirow{3}{*}{\makecell{G2-G3}} & Diagonal & C & 3120 &  & 250 ---  \\[6pt]
\cline{2-6}
 & Top chord & C & 2500 &  & 250 ---  \\[6pt]
\cline{2-6}
 & Bottom chord & T & 2500 &  & 400 ---  \\[6pt]
\hline
\multirow{3}{*}{\makecell{G3-G4}} & Diagonal & C & 3120 &  & 250 ---  \\[6pt]
\cline{2-6}
 & Top chord & C & 2500 &  & 250 ---  \\[6pt]
\cline{2-6}
 & Bottom chord & T & 2500 &  & 400 ---  \\[6pt]
\hline
\end{longtable}
\noindent\textit{Note:  3. Limit = 250 for compression members, 400 for tension members. $K = 1.0$ for members with both ends pinned.}

\vspace{1em}
\begin{longtable}{|C{2.0cm}|L{1.8cm}|C{2.2cm}|C{3.0cm}|C{1.8cm}|C{1.8cm}|C{1.2cm}|C{1.8cm}|}
\caption{\textbf{Cross Bracing Design --- Capacity Summary}}\label{tab:cb-capacity}\\
\hline
\textbf{Panel} & \textbf{Member} & \textbf{Section} & \textbf{Governing LC} & \textbf{Demand (kN)} & \textbf{Capacity (kN)} & \textbf{UR} & \textbf{Status} \\
\hline
\endfirsthead

\multicolumn{8}{c}{\small\textit{\tablename\ \thetable\ -- Continued from previous page}} \\
\hline
\textbf{Panel} & \textbf{Member} & \textbf{Section} & \textbf{Governing LC} & \textbf{Demand (kN)} & \textbf{Capacity (kN)} & \textbf{UR} & \textbf{Status} \\
\hline
\endhead

\hline
\multicolumn{8}{r}{\footnotesize\textit{Continued on next page\ldots}} \\
\endfoot

\hline
\endlastfoot
\multirow{2}{*}{\makecell{G1-G2}} & Diagonal &  & LC-ULS-4 & -168.500 &  &  &  \\[6pt]\cline{2-8}
 & Chord &  &  & -92.400 &  &  &  \\[6pt]\cline{2-8}
\hline
\multirow{2}{*}{\makecell{G2-G3}} & Diagonal &  & LC-ULS-4 & -152.000 &  &  &  \\[6pt]\cline{2-8}
 & Chord &  &  & -84.200 &  &  &  \\[6pt]\cline{2-8}
\hline
\multirow{2}{*}{\makecell{G3-G4}} & Diagonal &  & LC-ULS-4 & -168.500 &  &  &  \\[6pt]\cline{2-8}
 & Chord &  &  & -92.400 &  &  &  \\[6pt]\cline{2-8}
\hline
\end{longtable}
\noindent\textit{Note: Designed per IS 800 Cl. 7 (compression) and Cl. 6 (tension). OsdagBridge cross-bracing module used.}

% ===========================
\section{End Diaphragm Design}
\label{sec:end-diaphragm}
% ===========================

End diaphragms at the supports transfer transverse loads to the bearings, restrain the bottom flanges against lateral displacement, and maintain the girder cross-section geometry during construction and in service. They are designed per IS~800:2007 and IRC~24:2010 Cl.~507.

\vspace{1em}

\vspace{0.4em}
\noindent
\setlength{\tabcolsep}{4pt}
\setlength{\LTleft}{0pt}
\setlength{\LTright}{\fill}

\begin{longtable}{|C{2.0cm}|L{2.0cm}|L{2.2cm}|C{2.5cm}|C{2.0cm}|C{2.0cm}|}
\caption{\textbf{End Diaphragm --- Connection and Section Properties}}\label{tab:ed-section-props}\\
\hline
\textbf{Panel} & \textbf{Member} & \textbf{Connection} & \textbf{Section} & \textbf{$A_g$ (mm$^2$)} & \textbf{$r_{min}$ (mm)} \\
\hline
\endfirsthead

\multicolumn{6}{c}{\small\textit{\tablename\ \thetable\ -- Continued from previous page}} \\
\hline
\textbf{Panel} & \textbf{Member} & \textbf{Connection} & \textbf{Section} & \textbf{$A_g$ (mm$^2$)} & \textbf{$r_{min}$ (mm)} \\
\hline
\endhead

\hline
\multicolumn{6}{r}{\footnotesize\textit{Continued on next page\ldots}} \\
\endfoot

\hline
\endlastfoot
\multirow{2}{*}{\makecell{G1-G2}} & Diagonal & Bolted &  & 1550.0 &  \\[6pt]\cline{2-6}
\multirow{2}{*}{\makecell{G1-G2}} & Top / Bottom chord & Bolted &  &  &  \\[6pt]\cline{2-6}
\hline
\multirow{2}{*}{\makecell{G2-G3}} & Diagonal &  &  &  &  \\[6pt]\cline{2-6}
\multirow{2}{*}{\makecell{G2-G3}} & Top / Bottom chord &  &  &  &  \\[6pt]\cline{2-6}
\hline
\multirow{2}{*}{\makecell{G3-G4}} & Diagonal &  &  &  &  \\[6pt]\cline{2-6}
\multirow{2}{*}{\makecell{G3-G4}} & Top / Bottom chord &  &  &  &  \\[6pt]\cline{2-6}
\hline
\end{longtable}
\noindent\textit{Note: $A_g$ = gross cross-sectional area; $r_{min}$ = minimum radius of gyration.}

\vspace{1em}
\begin{longtable}{|C{2.2cm}|L{2.2cm}|L{2.5cm}|C{2.5cm}|C{2.5cm}|>{\centering\arraybackslash}p{3.6cm}|}
\caption{\textbf{End Diaphragm --- Slenderness Ratio Check (IS~800 Cl.~3.8)}}\label{tab:ed-slenderness}\\
\hline
\textbf{Panel} & \textbf{Member} & \textbf{Nature} & \textbf{Eff.\ Length $KL$ (mm)} & \textbf{$KL/r$} & \textbf{Limit / Status} \\
\hline
\endfirsthead

\multicolumn{6}{c}{\small\textit{\tablename\ \thetable\ -- Continued from previous page}} \\
\hline
\textbf{Panel} & \textbf{Member} & \textbf{Nature} & \textbf{Eff.\ Length $KL$ (mm)} & \textbf{$KL/r$} & \textbf{Limit / Status} \\
\hline
\endhead

\hline
\multicolumn{6}{r}{\footnotesize\textit{Continued on next page\ldots}} \\
\endfoot

\hline
\endlastfoot
\multirow{3}{*}{\makecell{G1-G2}} & Diagonal & C & 3120 &  & 250 ---  \\[6pt]
\cline{2-6}
 & Top chord & C & 2500 &  & 250 ---  \\[6pt]
\cline{2-6}
 & Bottom chord & T & 2500 &  & 400 ---  \\[6pt]
\hline
\multirow{3}{*}{\makecell{G2-G3}} & Diagonal & C & 3120 &  & 250 ---  \\[6pt]
\cline{2-6}
 & Top chord & C & 2500 &  & 250 ---  \\[6pt]
\cline{2-6}
 & Bottom chord & T & 2500 &  & 400 ---  \\[6pt]
\hline
\multirow{3}{*}{\makecell{G3-G4}} & Diagonal & C & 3120 &  & 250 ---  \\[6pt]
\cline{2-6}
 & Top chord & C & 2500 &  & 250 ---  \\[6pt]
\cline{2-6}
 & Bottom chord & T & 2500 &  & 400 ---  \\[6pt]
\hline
\end{longtable}
\noindent\textit{Note:  3. Limit = 250 for compression members, 400 for tension members. $K = 1.0$ for members with both ends pinned.}

\vspace{1em}

\begin{longtable}{|C{2.0cm}|L{1.8cm}|C{2.2cm}|C{3.0cm}|C{1.8cm}|C{1.8cm}|C{1.2cm}|C{1.8cm}|}
\caption{\textbf{End Diaphragm Design --- Capacity Summary}}\label{tab:ed-capacity}\\
\hline
\textbf{Panel} & \textbf{Member} & \textbf{Section} & \textbf{Governing LC} & \textbf{Demand (kN)} & \textbf{Capacity (kN)} & \textbf{UR} & \textbf{Status} \\
\hline
\endfirsthead

\multicolumn{8}{c}{\small\textit{\tablename\ \thetable\ -- Continued from previous page}} \\
\hline
\textbf{Panel} & \textbf{Member} & \textbf{Section} & \textbf{Governing LC} & \textbf{Demand (kN)} & \textbf{Capacity (kN)} & \textbf{UR} & \textbf{Status} \\
\hline
\endhead

\hline
\multicolumn{8}{r}{\footnotesize\textit{Continued on next page\ldots}} \\
\endfoot

\hline
\endlastfoot
\multirow{2}{*}{\makecell{G1-G2}} & Diagonal &  & LC-ULS-4 & -168.500 &  &  &  \\[6pt]\cline{2-8}
 & Chord &  &  & -92.400 &  &  &  \\[6pt]\cline{2-8}
\hline
\multirow{2}{*}{\makecell{G2-G3}} & Diagonal &  & LC-ULS-4 & -152.000 &  &  &  \\[6pt]\cline{2-8}
 & Chord &  &  & -84.200 &  &  &  \\[6pt]\cline{2-8}
\hline
\multirow{2}{*}{\makecell{G3-G4}} & Diagonal &  & LC-ULS-4 & -168.500 &  &  &  \\[6pt]\cline{2-8}
 & Chord &  &  & -92.400 &  &  &  \\[6pt]\cline{2-8}
\hline
\end{longtable}
\noindent\textit{Note: Designed per IS 800 Cl. 7 (compression) and Cl. 6 (tension). OsdagBridge cross-bracing module used.}

% ===========================
\section{Overall Design Check Summary}
\label{sec:overall-summary}
% ===========================


\begin{figure}[H]
\centering
\includegraphics[width=0.88\textwidth]{C:/Users/Akshat/AppData/Local/Temp/tmp27f6l637/images/ur_summary_chart.png}
\caption{Overall Design Check Summary --- Component Utilization Ratios (UR = Demand / Capacity)}
\label{fig:ur-summary-chart}
\end{figure}


\vspace{0.8em}
\begin{longtable}{|C{3.4cm}|L{4.5cm}|C{2.3cm}|C{2.3cm}|>{\centering\arraybackslash}p{1.6cm}|}
\caption{\textbf{Overall Design Check Summary --- All Members}}\label{tab:overall-summary}\\
\hline
\textbf{Member / Check} & \textbf{Governing Load Combo} & \textbf{Demand} & \textbf{Capacity} & \textbf{UR} \\
\hline
\endfirsthead

\multicolumn{5}{c}{\small\textit{\tablename\ \thetable\ -- Continued from previous page}} \\
\hline
\textbf{Member / Check} & \textbf{Governing Load Combo} & \textbf{Demand} & \textbf{Capacity} & \textbf{UR} \\
\hline
\endhead

\hline
\multicolumn{5}{r}{\footnotesize\textit{Continued on next page\ldots}} \\
\endfoot

\hline
\endlastfoot

Girder --- Moment & --- & --- & --- & --- \\[6pt]
\hline
Girder --- Shear & --- & --- & --- & --- \\[6pt]
\hline
Girder --- LTB (constr.) & --- & --- & --- & --- \\[6pt]
\hline
Girder --- Deflection & --- & --- & --- & --- \\[6pt]
\hline
Girder --- Stress & --- & --- & --- & --- \\[6pt]
\hline
Girder --- Fatigue & --- & --- & --- & --- \\[6pt]
\hline
Transverse Shear (slab) & --- & --- & --- & --- \\[6pt]
\hline
Crack Width (slab) & Frequent SLS & 0.180 mm & 0.250 mm & 0.72 \\[6pt]
\hline
Deck --- Flexure (sagging) & Basic ULS: 0DL + 0LL & 145.50 kN-m/m & 185.00 kN-m/m & 0.79 \\[6pt]
\hline
Deck --- Flexure (hogging) & Basic ULS: 0DL + 0LL & 120.00 kN-m/m & 160.00 kN-m/m & 0.75 \\[6pt]
\hline
Deck --- Cantilever Overhang & Basic ULS: 0DL + 0LL & N/A & N/A & N/A \\[6pt]
\hline
Deck --- Punching Shear & Basic ULS: 0DL + 0LL & 245.00 MPa & 420.00 MPa & 0.58 \\[6pt]
\hline
Deck --- One-Way Shear & Basic ULS: 0DL + 0LL & 165.00 kN/m & 240.00 kN/m & 0.69 \\[6pt]
\hline
Cross Bracing --- Compression & --- & --- & --- & --- \\[6pt]
\hline
Cross Bracing --- Tension & --- & --- & --- & --- \\[6pt]
\hline
Cross Bracing --- Slenderness & --- & --- & --- & --- \\[6pt]
\hline
End Diaphragm --- Moment & \multicolumn{4}{c|}{Rolled / Welded section --- design to be added} \\[6pt]
\hline
End Diaphragm --- Shear & \multicolumn{4}{c|}{Rolled / Welded section --- design to be added} \\[6pt]
\hline
\end{longtable}
\noindent\textit{Note: UR = Demand / Capacity. All values $\leq 1.0$ indicate passing checks. The governing check for each component is highlighted in the individual design check sections above.}



\chapter{Drawings and Visualizations}
\label{ch:drawings}

This section presents CAD-generated views of the designed bridge and its components. All views are generated automatically by OsdagBridge using pythonOCC.

\section{Bridge Configuration and Layout}
\label{sec:bridge-layout}

\begin{figure}[H]
\centering
\fbox{\parbox{0.97\textwidth}{
\textit{[ PLACEHOLDER:  Overall 3D Bridge Superstructure ]}
}}
\caption{Overall 3D Bridge Superstructure}
\label{fig:bridge-3d}
\end{figure}


\begin{figure}[H]
\centering
\fbox{\parbox{0.97\textwidth}{
\textit{[ PLACEHOLDER:  Typical Cross Section ]}
}}
\caption{Typical Cross Section}
\label{fig:cross-section}
\end{figure}


\begin{figure}[H]
\centering
\fbox{\parbox{0.97\textwidth}{
\textit{[ PLACEHOLDER:  Top View ]}
}}
\caption{Top View}
\label{fig:top-view}
\end{figure}


\section{Plate Girder --- Detailed Views}
\label{sec:girder-views}

\begin{figure}[H]
\centering
\fbox{\parbox{0.97\textwidth}{
\textit{[ PLACEHOLDER:  3D View of Plate Girders ]}
}}
\caption{3D View of Plate Girders}
\label{fig:girder-3d}
\end{figure}


\begin{figure}[H]
\centering
\fbox{\parbox{0.97\textwidth}{
\textit{[ PLACEHOLDER:  Cross Section of Plate Girder ]}
}}
\caption{Cross Section of Plate Girder}
\label{fig:girder-xsec}
\end{figure}


\begin{figure}[H]
\centering
\fbox{\parbox{0.97\textwidth}{
\textit{[ PLACEHOLDER:  Side View of Girder ]}
}}
\caption{Side View of Girder}
\label{fig:girder-side}
\end{figure}


\section{Cross Bracing Detail}
\label{sec:bracing-detail}

\begin{figure}[H]
\centering
\fbox{\parbox{0.97\textwidth}{
\textit{[ PLACEHOLDER:  Cross Bracing Layout ]}
}}
\caption{Cross Bracing Layout}
\label{fig:cb-layout}
\end{figure}


\noindent
\begin{minipage}[t]{0.31\textwidth}
\centering
\fbox{\parbox{0.95\linewidth}{\centering\textit{[ Cross Bracing -- Bracing Section ]}}}
\captionof{figure}{Cross Bracing -- Bracing Section}
\label{fig:cb-bracing-sec}
\end{minipage}
\hfill
\begin{minipage}[t]{0.31\textwidth}
\centering
\fbox{\parbox{0.95\linewidth}{\centering\textit{[ Cross Bracing -- Top Chord Section ]}}}
\captionof{figure}{Cross Bracing -- Top Chord Section}
\label{fig:cb-top-chord}
\end{minipage}
\hfill
\begin{minipage}[t]{0.31\textwidth}
\centering
\fbox{\parbox{0.95\linewidth}{\centering\textit{[ Cross Bracing -- Bottom Chord Section ]}}}
\captionof{figure}{Cross Bracing -- Bottom Chord Section}
\label{fig:cb-bottom-chord}
\end{minipage}
            

\section{End Diaphragm Detail}
\label{sec:diaphragm-detail}

\begin{figure}[H]
\centering
\fbox{\parbox{0.97\textwidth}{
\textit{[ PLACEHOLDER:  End Diaphragm Layout ]}
}}
\caption{End Diaphragm Layout}
\label{fig:ed-layout}
\end{figure}


\noindent
\begin{minipage}[t]{0.31\textwidth}
\centering
\fbox{\parbox{0.95\linewidth}{\centering\textit{[ End Diaphragm -- Bracing Section ]}}}
\captionof{figure}{End Diaphragm -- Bracing Section}
\label{fig:ed-bracing-sec}
\end{minipage}
\hfill
\begin{minipage}[t]{0.31\textwidth}
\centering
\fbox{\parbox{0.95\linewidth}{\centering\textit{[ End Diaphragm -- Top Chord Section ]}}}
\captionof{figure}{End Diaphragm -- Top Chord Section}
\label{fig:ed-top-chord}
\end{minipage}
\hfill
\begin{minipage}[t]{0.31\textwidth}
\centering
\fbox{\parbox{0.95\linewidth}{\centering\textit{[ End Diaphragm -- Bottom Chord Section ]}}}
\captionof{figure}{End Diaphragm -- Bottom Chord Section}
\label{fig:ed-bottom-chord}
\end{minipage}



\chapter{Material Take-off \& Quantity Summary}
\label{ch:material-takeoff}

This section summarizes the material quantities calculated for the superstructure, including structural steel plate girders, cross bracing, end diaphragms, concrete deck slab, reinforcement steel, and shear connectors.

\vspace{0.8em}
\begingroup
\setlength{\tabcolsep}{3.5pt}
\begin{longtable}{|C{1.0cm}|L{3.8cm}|C{2.6cm}|C{1.8cm}|C{1.8cm}|C{1.8cm}|C{1.8cm}|}
\caption{\textbf{Bill of Materials (Structural Steel, Concrete, and Reinforcement Quantities)}}\label{tab:bill-of-materials}\\
\hline
\textbf{S.N.} & \textbf{Item Description} & \textbf{Volume} & \textbf{Quantity} & \textbf{Total Volume} & \textbf{Weight (MT)} & \textbf{Total Weight (MT)} \\
\hline
\endfirsthead

\multicolumn{7}{c}{\small\textit{\tablename\ \thetable\ -- Continued from previous page}} \\
\hline
\textbf{S.N.} & \textbf{Item Description} & \textbf{Volume} & \textbf{Quantity} & \textbf{Total Volume} & \textbf{Weight (MT)} & \textbf{Total Weight (MT)} \\
\hline
\endhead

\hline
\multicolumn{7}{r}{\footnotesize\textit{Continued on next page\ldots}} \\
\endfoot

\hline
\endlastfoot

1 & Structural Steel (IS 2062) for Girders & $0.06260\text{ m}^2 \times 30.00\text{ m} = 1.87800\text{ m}^3$ & 4 & 7.51 & 14.74 & 58.97 \\
\hline
2(a) & Cross Bracing - Top Chord & $0.00155\text{ m}^2 \times 2.50\text{ m} = 0.00387\text{ m}^3$ & 15 & 0.06 & 0.0304 & 0.46 \\
\hline
2(b) & Cross Bracing - Bottom Chord & $0.00155\text{ m}^2 \times 2.50\text{ m} = 0.00387\text{ m}^3$ & 15 & 0.06 & 0.0304 & 0.46 \\
\hline
2(c) & Cross Bracing - Diagonal Chord & $0.00155\text{ m}^2 \times 3.12\text{ m} = 0.00484\text{ m}^3$ & 30 & 0.15 & 0.0380 & 1.14 \\
\hline
3 & Concrete (M40) for Deck Slab & $10.50\text{ m} \times 0.25\text{ m} \times 30.00\text{ m} = 78.75\text{ m}^3$ & 1 & 78.75 & 196.88 & 196.88 \\
\hline
4 & Reinforcement Steel (Fe 500) & $0.040127\text{ m}^2 \times 30.00\text{ m} = 1.20382\text{ m}^3$ & 1 & 1.20 & 9.45 & 9.45 \\
\hline
5 & Shear Stud Connectors & $0.000380\text{ m}^2 \times 0.150\text{ m} = 0.000057\text{ m}^3$ & 2412 & 0.14 & 0.000448 & 1.080 \\
\hline
6 & Crash Barrier & N.A. & N.A. & N.A. & N.A. & N.A. \\
\hline
\end{longtable}
\endgroup

\section{Material Quantity Charts}
\label{sec:material-charts}


\begin{figure}[H]
\centering
\includegraphics[width=0.88\textwidth]{C:/Users/Akshat/AppData/Local/Temp/tmp27f6l637/images/material_steel_tonnage.png}
\caption{Structural Steel Tonnage Breakdown by Component (IS 2062)}
\label{fig:chart-steel-tonnage}
\end{figure}

\vspace{0.8em}

\begin{figure}[H]
\centering
\includegraphics[width=0.88\textwidth]{C:/Users/Akshat/AppData/Local/Temp/tmp27f6l637/images/material_concrete_rebar.png}
\caption{Concrete Volume and Reinforcement Steel Weight Take-Off}
\label{fig:chart-concrete-rebar}
\end{figure}



\chapter{Standards \& Assumptions}
\label{ch:Design Standards}

This section provides references to standards used to calibrate the OsdagBridge
design modules, and notes the limitations of the current software version.

\section{Design Standards}
\label{sec:design_standards}

The following Indian Road Congress (IRC) codes and Indian Standards (IS) 
form the basis of all design calculations in this software.

\vspace{0.5cm}

\begin{longtable}{|C{2.5cm}|C{2.0cm}|L{12.0cm}|}
\caption{\textbf{IRC Codes}}\label{tab:irc-codes}\\
\hline
\textbf{Code} & \textbf{Year} & \textbf{Title / Scope} \\
\hline
\endfirsthead

\multicolumn{3}{c}{\small\textit{\tablename\ \thetable\ -- Continued from previous page}} \\
\hline
\textbf{Code} & \textbf{Year} & \textbf{Title / Scope} \\
\hline
\endhead

\hline
\multicolumn{3}{r}{\footnotesize\textit{Continued on next page\ldots}} \\
\endfoot

\hline
\endlastfoot

IRC 5 & 2015 & General Features of Design - carriageway widths, kerb, footpath dimensions \\ 
\hline
IRC 6 & 2017 & Loads and Load Combinations - dead load, live load, impact, wind, temperature, etc. \\ 
\hline
IRC 22 & 2015 & Composite Construction (LS) - Composite section properties, ULS/SLS design, shear connectors \\ 
\hline
IRC 24 & 2010 & Steel Road Bridges (LS) - Stiffener design, skew angle limits, diaphragm requirements \\ 
\hline
IRC 112 & 2020 & Concrete Road Bridges - Deck slab flexure, shear, crack width, reinforcement \\ 
\hline
IRC SP 114 & 2018 & Seismic Design of Road Bridges \\ 
\hline
\end{longtable}

\begin{longtable}{|C{2.5cm}|C{2.0cm}|L{12.0cm}|}
\caption{\textbf{IS Codes}}\label{tab:is-codes}\\
\hline
\textbf{Code} & \textbf{Year} & \textbf{Scope} \\
\hline
\endfirsthead

\multicolumn{3}{c}{\small\textit{\tablename\ \thetable\ -- Continued from previous page}} \\
\hline
\textbf{Code} & \textbf{Year} & \textbf{Scope} \\
\hline
\endhead

\hline
\multicolumn{3}{r}{\footnotesize\textit{Continued on next page\ldots}} \\
\endfoot

\hline
\endlastfoot

IS 800 & 2007 & Steel construction - tension, compression, bending, shear, LTB, stiffeners, combined checks \\
\hline
IS 456 & 2000 & Concrete - simplified stress-block for deck moment capacity \\
\hline
IS 1786 & 2008 & Reinforcement steel properties \\
\hline
IS 1893 (Part 3) & 2014 & Earthquake resistant design \\
\hline
IS 2062 & 2011 & Structural steel - yield and ultimate strength by grade \\
\hline
\end{longtable}

\clearpage
\section{Analysis and Design Assumptions of This Version}
\label{sec:assumptions}

\textbf{Structural Analysis}

\begin{itemize}
    \item All girders are modelled as simply supported; continuous spans are not currently supported.
    \item A 3D grillage model (OSPGrillage) is used for load distribution. Grillage members carry composite section properties after the construction stage.
    \item The transverse member forces are computed using an approximate 2D frame analogy. For irregular or skewed geometries, a 3D FEM is recommended.
    \item Fixed bearing stiffness is modelled as $k = 1{,}000{,}000 \,\text{kN/m}$ (virtually rigid); free/expansion bearing as $k = 100 \, \text{kN/m}$.
    \item Construction stage sequence analysis is approximate. Detailed staged analysis should be performed for long-term deflection checks.
\end{itemize}

\textbf{Composite Action}

\begin{itemize}
    \item Full shear connection is assumed at ULS with headed stud connectors designed per IRC 22:2015 Cl.606.
    \item Short-term composite section properties (modular ratio $n = E_s/E_{cm}$) are used for ULS checks.
    \item Long-term composite section properties (with creep-adjusted modular ratio) are used for SLS deflection and crack-width checks.
    \item Pre-composite stage: the steel girder alone resists all construction loads prior to concrete gaining strength.
\end{itemize}


\textbf{Material Properties (IRC 22:2015 Annex III)}

\begin{itemize}
\item Steel: $E_s = 200{,}000 \, \text{MPa}, \ G_s = 80{,}000 \, \text{MPa}, \ \nu = 0.30, \ \alpha = 11.7 \times 10^{-6}/^\circ\mathrm{C}$ (grade-independent).
\item Minimum structural concrete grade: M25.
\item Default reinforcement grade: Fe500.
\end{itemize}

\textbf{Partial Safety Factors (IRC 22:2015 Cl.601.4)}

\begin{itemize}
\item $\gamma_{m0}$ (steel yield, ULS) = 1.10
\item $\gamma_{m1}$ (steel ultimate, ULS) = 1.25
\item Reinforcement (ULS) = 1.15
\item Welds -- shop: 1.25; field: 1.50
\item Fatigue ($\gamma_{mft}$) = 1.35
\end{itemize}

\textbf{Loading}

\begin{itemize}
\item Dead load densities per IRC 6:2017 Cl.203: structural steel $78.5 \ \text{kN/m}^3$, concrete $25.0 \ \text{kN/m}^3$, bituminous wearing course $24.0 \ \text{kN/m}^3$.
\item Multi-lane live load reduction factors per IRC 6:2017 Cl.204.4 Table 6A: 1st lane = 1.0, 2nd lane = 0.8, 3rd lane onwards = 0.4.
\item Impact factor (dynamic load allowance) computed from span per IRC 6:2017 Cl.208.2/208.3.
\item Wind load applied as transverse, longitudinal, and vertical components per IRC 6:2017 Cl.209.3.3--209.3.5.
\end{itemize}

\textbf{Serviceability Limits}

\begin{itemize}
\item Deflection limits (IRC 22:2015 Cl.604.3.2): live load + impact $\leq L/800$; total $\leq L/600$.
\item SLS stress limits: concrete $\sigma_c \leq 0.48 f_{ck}$; rebar $\sigma_s \leq 0.80 f_{yk}$; steel $f_e \leq 0.9 f_y$.
\item Permissible crack width: $w_k \leq 0.3 \ \text{mm}$ (bridge deck, exposure class XS2/XD2 per IRC 112:2020 Cl.12.3.2).
\end{itemize}

\textbf{Fatigue}

\begin{itemize}
\item Reference fatigue life: $N_{sc} = 2 \times 10^6$ cycles (IRC 22:2015 Cl.605).
\item Constant stress range is assumed. A thickness correction factor $\mu_r$ is applied for plate thickness $> 25 \ \text{mm}$.
\end{itemize}

\textbf{Stiffener Design}

\begin{itemize}
\item Intermediate transverse stiffener and bearing stiffener design follows IS 800:2007 Cl.8.7.2/8.7.3 and IRC 24:2010 Cl.509.7.2/509.7.3.
\end{itemize}


\section{Known Limitations of This Version}
\label{sec:limitations}

\begin{itemize}
\item Substructure (piers, pile caps, foundations) and bearing design are not included.
\item Splice connection design is not implemented.
\item Skew angle $>$ 15 degrees requires independent manual analysis (IRC 24 Cl. 504.8).
\item Construction stage sequence analysis is approximate; detailed staged analysis
  should be performed for long-term deflection checks.
\item The grillage analysis assumes simply supported boundary conditions;
  continuous spans are not currently supported.
\end{itemize}


\chapter*{References}
\addcontentsline{toc}{chapter}{References}


\begin{enumerate}

\item IRC 5 (2024) --- \textit{Standard Specifications and Code of Practice for Road Bridges, Section I: General Features of Design.}

\item IRC 6 (2017) --- \textit{Standard Specifications and Code of Practice for Road Bridges, Section II: Loads and Load Combinations.}

\item IRC 22 (2014) --- \textit{Standard Specifications and Code of Practice for Road Bridges, Section VI: Composite Construction (Limit State Design).}

\item IRC 24 (2010) --- \textit{Standard Specifications and Code of Practice for Road Bridges, Section V: Steel Road Bridges (Limit State Method).}

\item IRC 112 (2011) --- \textit{Code of Practice for Concrete Road Bridges.}

\item IRC SP 114 (2018) --- \textit{Guidelines for Seismic Design of Road Bridges.}

\item IS 800 (2007) --- \textit{Indian Standard: General Construction in Steel --- Code of Practice.}

\item IS 2062 (2011) --- \textit{Hot Rolled Medium and High Tensile Structural Steel --- Specification.}

\item Subramanian, N. (2008). \textit{Design of Steel Structures.} Oxford University Press.

\item Steel-INSDAG Teaching Resource Materials. \url{https://www.steel-insdag.org}

\end{enumerate}

\end{document}